% concatenated from main.tex
\pdfoutput=1 % ensure arXiv direct PDF output

\documentclass[11pt]{amsart}

\usepackage{amssymb,amsthm,enumitem,colonequals,tikz-cd,microtype}
\usepackage[normalem]{ulem} 

%% Font style: old versions mlmodern, \usepackage[cal=euler,bfcal,bb=px,bfbb]{mathalpha}
\usepackage[osf]{Baskervaldx}
\usepackage[baskervaldx]{newtxmath}
\usepackage[cal=boondoxo]{mathalfa} % mathcal from STIX, unslanted a bit

%% Geometry with custom margins
\usepackage[top=3.75cm, bottom=3cm, left=3.5cm, right=3.5cm]{geometry}

%% High penalty to avoid footnote breaks
\interfootnotelinepenalty=10000 

%% Color definitions for links and references
\usepackage{xcolor}
\colorlet{darkblue}{blue!55!black}
\colorlet{darkcyan}{cyan!50!black}
\colorlet{darkgreen}{green!60!black}

%% URL hyphenation option must come before hyperref
\PassOptionsToPackage{hyphens}{url}
\usepackage{hyperref}
\hypersetup{
    colorlinks=true,
    linkcolor=darkblue,
    urlcolor=darkcyan,
    citecolor=darkgreen,
}

%% Colored eqref for emphasis (optional)
\def\eqref#1{\textcolor{darkblue}{(\ref{#1})}}

%% Cleveref utility
\usepackage[nameinlink]{cleveref} 
\Crefformat{section}{#2\S#1#3}
\Crefmultiformat{section}{#2\S\S#1#3}{ and~#2#1#3}{, #2#1#3}{, and~#2#1#3}

%% Overfill notice
\usepackage[pagewise]{lineno}
\overfullrule=100pt % for debugging, comment or set to 0 for final

\let\oldequation\equation
\let\oldendequation\endequation
\renewenvironment{equation}{\linenomathNonumbers\oldequation}{\oldendequation\endlinenomath}
\expandafter\let\expandafter\oldequationstar\csname equation*\endcsname
\expandafter\let\expandafter\oldendequationstar\csname endequation*\endcsname
\renewenvironment{equation*}{\linenomathNonumbers\oldequationstar}{\oldendequationstar\endlinenomath}
\let\oldalign\align
\let\oldendalign\endalign
\renewenvironment{align}{\linenomathNonumbers\oldalign}{\oldendalign\endlinenomath}
\expandafter\let\expandafter\oldalignstar\csname align*\endcsname
\expandafter\let\expandafter\oldendalignstar\csname endalign*\endcsname
\renewenvironment{align*}{\linenomathNonumbers\oldalignstar}{\oldendalignstar\endlinenomath}

%% Theorem and counter setup
\theoremstyle{plain}
\newtheorem{theorem}{Theorem}[section]
\newtheorem{lemma}[theorem]{Lemma}
\newtheorem{corollary}[theorem]{Corollary}
\newtheorem{proposition}[theorem]{Proposition}

\theoremstyle{definition}
\newtheorem{definition}[theorem]{Definition}
\newtheorem{example}[theorem]{Example}
\newtheorem{remark}[theorem]{Remark}
\newtheorem{setup}[theorem]{Setup}
\newtheorem{convention}[theorem]{Convention}

\newtheorem*{disclaimer}{Disclaimer}
\newtheorem*{ack}{Acknowledgments}

% declare alias for cleveref, otherwise it doesn't do what it should; this probably isn't best fix but found it on the web
% \AddToHook{env/theorem/begin}{\crefalias{intro}{theorem}}
\AddToHook{env/conjecture/begin}{\crefalias{theorem}{conjecture}}
\AddToHook{env/lemma/begin}{\crefalias{theorem}{lemma}}
\AddToHook{env/corollary/begin}{\crefalias{theorem}{corollary}}
\AddToHook{env/proposition/begin}{\crefalias{theorem}{proposition}}
\AddToHook{env/definition/begin}{\crefalias{theorem}{definition}}
\AddToHook{env/remark/begin}{\crefalias{theorem}{remark}}
\AddToHook{env/setup/begin}{\crefalias{theorem}{setup}}
\AddToHook{env/example/begin}{\crefalias{theorem}{example}}
\AddToHook{env/convention/begin}{\crefalias{theorem}{convention}}

%% toc and numbering depths
\setcounter{tocdepth}{1}
\setcounter{secnumdepth}{4}
\numberwithin{equation}{section}
\numberwithin{theorem}{section}

%% Comment out collaborator notes before final version
\newcommand{\Pat}[1]{{\color{blue} $\clubsuit$ Pat: [#1]}}
\newcommand{\Michal}[1]{{\color{teal} $\heartsuit$ Michal: [#1]}}
\newcommand{\Simone}[1]{{\color{orange} $\spadesuit$ Simone: [#1]}}


%%% DOCUMENT INFORMATION --------------------------------------------------
\title[Perfectly generated $t$-structures, algebraic spaces]{Perfectly generated $t$-structures \\ for algebraic spaces}

\author[M.~Hrbek]{Michal Hrbek}
\address{M.~Hrbek,
Institute of Mathematics of the Czech Academy of Sciences,
Žitná 25, 115 67 Prague, Czechia}
\email{hrbek@math.cas.cz}

\author[P.~Lank]{Pat Lank}
\address{P.~Lank,
Dipartimento di Matematica “F. Enriques”, Universit\`{a} degli Studi di Milano, Via Cesare
Saldini 50, 20133 Milano, Italy}
\email{plankmathematics@gmail.com}

\author[S.~Pizzirani]{Simone Pizzirani}
\address{S.~Pizzirani, \\
Institute of Mathematics of the Czech Academy of Sciences,
Žitná 25, 115 67 Prague, Czechia and Faculty of Mathematics and Physics, Department
of Algebra, Charles University, Sokolovsk\'a 83, 186 75 Prague, Czechia}
\email{pizzirani@math.cas.cz}

\date{\today}

\keywords{Derived categories, algebraic spaces, $t$-structures, classification, perfect complexes}

\subjclass[2020]{14A30 (primary), 14F08, 18G80, 14D23}

% 14A30 Fundamental constructions in algebraic geometry involving higher and derived categories (homotopical algebraic geometry, derived algebraic geometry, etc.) 

% 14D23 Stacks and moduli problems

% 14F08 Derived categories of sheaves, dg categories, and related constructions in algebraic geometry

% 18G80 Derived categories, triangulated categories

\begin{document}
    
\begin{abstract}
    This work studies $t$-structures for the derived category of complexes with quasi-coherent cohomology on algebraic spaces. 
    The main result shows that the relative standard $t$-structures are compactly generated. 
    As an application, we classify compactly generated tensor $t$-structures via Thomason filtrations on the underlying topological space. 
\end{abstract}

\maketitle

\tableofcontents

%%%%%%%%%%%%%%%%%%%%%%%%%%%%%%%%%%%%%
\section{Introduction}
\label{sec:intro}
%%%%%%%%%%%%%%%%%%%%%%%%%%%%%%%%%%%%%

%%%%%%%%%%%%%%%%%%%%%%%%%%%%%%%%%%%%%
\subsection{What is known}
\label{sec:intro_what_is_known}
%%%%%%%%%%%%%%%%%%%%%%%%%%%%%%%%%%%%%

Since their introduction in \cite{Beilinson/Berstein/Deligne/Gabber:2018}, $t$-structures have become a fundamental tool for extracting information from derived categories in algebraic geometry and commutative algebra. 
Recent applications include the proof of the Antieau--Gepner--Heller conjecture \cite{Antieau/Gepner/Heller:2019}. 
On a Noetherian scheme, it has been shown that the category of perfect complexes $\operatorname{Perf}$ admits a bounded $t$-structure if, and only if, the scheme is regular \cite{Smith:2022,Neeman:2024}.
 
Roughly speaking, a $t$-structure on a triangulated category $\mathcal{T}$ is a pair of subcategories $(\mathcal{T}^{\leq 0}, \mathcal{T}^{\geq 0})$ satisfying certain conditions. 
This includes an orthogonality constraint and being able to express objects as extensions of the pair. 
In fact, the $t$-structure is completely determined by its \textit{aisle} $\mathcal{T}^{\leq 0}$ \cite{Keller/Vossieck:1988}. 

% For suitable algebraic spaces, it is known that the standard aisle $D^{\leq 0}_{\operatorname{qc}}(X)$ is `equivalent' to a $t$-structure that is singly compactly generated. 
% Here, equivalent means in the sense of \cite[Definition 0.14]{Neeman:2025}. 
% See \cite[Theorem 3.2(iii)]{Neeman:2024} for the case of schemes, and \cite[Lemma A.8]{GuisadoVillaalgordo/Lank:2026} for suitable algebraic spaces.

% More recently, it was shown that the standard aisle need not be itself singly compactly generated for many algebraic spaces \cite{Bhaduri/DeDeyn/Hrbek/Lank/ManaliRahul:2025}. 
% On the other hand, for Noetherian schemes, the standard aisle is compactly generated by perfect complexes \cite[Proposition 2.3]{Herbera/Hrbek/LeGros:2025}. 
% Beyond this, however, little appears to be known, even for quasi-compact quasi-separated schemes.

% Some historical context is helpful.
The study of $t$-structures on schemes has been very active. 
In full generality, a complete classification of $t$-structures is impossible even for affine schemes \cite{Stanley:2010}. 
Consequently, much of the literature has focused on classes of $t$-structures satisfying additional conditions.

For affine Noetherian schemes, a classification of compactly generated t-structures was obtained in \cite{AlsonoTarrio/JeremiasLopez/Saorin:2010}. Subsequently, \cite{Dubey/Sahoo:2023} extended these results to Noetherian schemes under a tensor compatibility assumption. 
The Noetherian hypothesis was later removed in the affine case by \cite{Hrbek:2020}, and those results play a key role in the present work.

These classifications are formulated in terms of `Thomason filtrations'. Namely, these are decreasing functions from the integers to the collection of Thomason subsets of the underlying topological space. 
See \Cref{sec:prelim_Thomason} for details. 
However, no analogous classification appears to be available for algebraic spaces, which is addressed in our work.

%%%%%%%%%%%%%%%%%%%%%%%%%%%%%%%%%%%%%
\subsection{What we do}
\label{sec:intro_what_we_do}
%%%%%%%%%%%%%%%%%%%%%%%%%%%%%%%%%%%%%

% %%%%%%%%%%%%%%%%%%%%%%%%%%%%%%%%%%%%%
% \subsubsection{Generation}
% \label{sec:intro_what_we_do_generation}
% %%%%%%%%%%%%%%%%%%%%%%%%%%%%%%%%%%%%%

Our first result is the following. 
It provides the key technical 
ingredient for our desired classification. 

\begin{proposition}
    \label{prop:psuedoapprox}
    Let $X$ be a quasi-compact quasi-separated algebraic space. Then $D^{\leq 0}_{\operatorname{qc},T}(X)$ is compactly generated for all closed subsets $T \subseteq |X|$ with quasi-compact complement.
\end{proposition}

\begin{corollary}
    \label{cor:psuedoapprox}
    Let $X$ be a quasi-compact quasi-separated scheme. Then $D^{\leq 0}_{\operatorname{qc},T}(X)$ is compactly generated for all closed subsets $T \subseteq |X|$ with quasi-compact complement.
\end{corollary}

\begin{remark}
    In the special case $T=|X|$, the corresponding $\infty$-categorical statement is \cite[Proposition 9.6.1.2]{Lurie:2018}, which asserts connective perfect generation for quasi-compact quasi-separated spectral algebraic spaces. 
    \Cref{prop:psuedoapprox} gives a triangulated refinement for quasi-compact quasi-separated algebraic spaces, and extends it the relative setting.
\end{remark}

To prove \Cref{prop:psuedoapprox}, we introduce a notion called \textit{pseudoapproximation}. 
This allows one to approximate bounded above complexes whose highest nonvanishing cohomology sheaf is finite type. 
The condition is weaker than approximation by perfect complexes \cite[Theorem 4.1]{Lipman/Neeman:2007} over Noetherian schemes. However, pseudoapproximation preserves cohomological supremum even in the general case, which is crucial for applications to $t$-structures. 
We show that quasi-compact quasi-separated schemes satisfy pseudoapproximation by perfect complexes. 
See \Cref{sec:pseudoapproximation} for details. 

It is not clear to us that quasi-compact quasi-separated algebraic spaces satisfy pseudoapproximation by perfect complexes.
In the Zariski topology, inductive arguments are available to glue along open covers. 
The situation is subtler for algebraic spaces.
Instead, we use elementary distinguished squares and the so-called induction principle on algebraic spaces to prove \Cref{prop:psuedoapprox}.
Key ingredients are \Cref{prop:psuedoapprox}, and tools developed to prove it, which allow us to glue complexes with prescribed cohomology and support.

We apply \Cref{prop:psuedoapprox} to obtain a classification of suitable $t$-structures on $D_{\operatorname{qc}}$. 
To be precise, an aisle $\mathcal{A}$ on $D_{\operatorname{qc}}(X)$ is called \textit{$\otimes$-closed} if $D^{\leq 0}_{\operatorname{qc}}(X)\otimes^{\mathbf{L}} \mathcal{A} \subseteq \mathcal{A}$. 
Applying \Cref{prop:psuedoapprox}, this is equivalent to tensoring by perfect complexes which belong to the standard aisle. 
See \Cref{rmk:proxy_aisle} for details.

Our main result is the following:

\begin{theorem}
    \label{thm:stacky_DB}
    Let $X$ be a quasi-compact quasi-separated algebraic space. 
    There is a one-to-one correspondence:
    \begin{displaymath}
        \begin{aligned}
            \{ \textrm{compactly generated } \otimes\!\textrm{\!-aisles on } & D_{\operatorname{qc}}(X)\} 
            \iff \{ \textrm{Thomason filtrations on } X\}.
        \end{aligned}
    \end{displaymath}
    In particular, for any $\otimes$-aisle on $D_{\operatorname{qc}}(X)$ generated by $\mathcal{P} \subseteq \operatorname{Perf}(X)$, we assign to it the Thomason filtration given by the rule (which is independent of the choice for $\mathcal{P}$)
    \begin{displaymath}
        n\mapsto \bigcup_{\substack{j\geq n\\ P\in \mathcal{P}}} \operatorname{supp}(\mathcal{H}^j (P)).
    \end{displaymath}
\end{theorem}

Notably, in the proof of \Cref{thm:stacky_DB}, our work allows us to bootstrap the affine results of \cite{Hrbek:2020} to go directly to algebraic spaces without any dependence on \cite{Dubey/Sahoo:2023}. 
To our knowledge, such a classification has not previously been established for arbitrary quasi-compact quasi-separated schemes.

% As a first step towards \Cref{thm:stacky_DB}, we show for an affine scheme that every Thomason filtration is supported by perfects. 
% Then we establish an auxiliary result (see \Cref{lem:pullback_t_exactness}). 
% It allows us to better understand such $t$-structures by passing to an affine scheme that is a smooth cover of the algebraic space. 
% This immediately enables us to bootstrap the classification known in the affine setting.

\begin{ack}
Hrbek and Pizzirani were supported by the project LQ100192601 Lumina quaeruntur, funded by the Czech Academy of Sciences (RVO 67985840). 
Lank was supported under the ERC Advanced Grant 101095900-TriCatApp, by the National Science Foundation under Grant No.\ DMS-1928930 while in residence at the Simons Laufer Mathematical Sciences Institute (formerly MSRI) in Spring 2024, and thanks the Institute of Mathematics of the Czech Academy of Sciences for its hospitality. 
The authors greatly appreciate discussions with Alexander Clark, Timothy De Deyn, Giovanna Le Gros, Kabeer Manali Rahul, Chris J.\ Parker, and Sergio Pavon. 
\end{ack}

%%%%%%%%%%%%%%%%%%%%%%%%%%%%%%%%%%%%%
\subsection*{Conventions/notation}
%%%%%%%%%%%%%%%%%%%%%%%%%%%%%%%%%%%%%

We follow \cite{StacksProject} for algebraic spaces and their derived categories. 
See \cite[\S 2.3]{GuisadoVillaalgordo/Lank:2026} for a summary.
Let $X$ be a quasi-compact quasi-separated algebraic space. 
$\operatorname{Mod}(X)$ is the Grothendieck abelian category of sheaves of $\mathcal{O}_X$-modules on the small \'{e}tale site $X_{\textrm{\'{e}tale}}$ of $X$ \cite[\href{https://stacks.math.columbia.edu/tag/03LT}{Tag 03LT}]{StacksProject}. 
$\operatorname{Qcoh}(X)$ is the full subcategory of $\operatorname{Mod}(X)$ consisting of quasi-coherent sheaves. 
$D(X) := D(\operatorname{Mod}(X))$ is the derived category of $\operatorname{Mod}(X)$. 
$D_{\operatorname{qc}}(X)$ is the full subcategory of $D(X)$ consisting of complexes with quasi-coherent cohomology sheaves. 
$\operatorname{Perf}(X)$ is the full subcategory of perfect complexes in $D_{\operatorname{qc}}(X)$. 
If $X$ is Noetherian, then $\operatorname{coh}(X)$ is the full subcategory of $\operatorname{Mod}(X)$ consisting of coherent sheaves and $D^b_{\operatorname{coh}}(X)$ denotes the full subcategory of $D(X)$ consisting of bounded pseudocoherent complexes.

%%%%%%%%%%%%%%%%%%%%%%%%%%%%%%%%%%%%%
\section{Preliminaries}
\label{sec:prelim}
%%%%%%%%%%%%%%%%%%%%%%%%%%%%%%%%%%%%%

We recall $t$-structures on triangulated categories, their tensor variants, and related topological constructions. 
Let $\mathcal{T}$ be a triangulated category admitting small coproducts. 
Its subcategory of compact objects is denoted by $\mathcal{T}^c$, consisting of those $c \in \mathcal{T}$ for which $\operatorname{Hom}(c,-)$ preserves coproducts.

%%%%%%%%%%%%%%%%%%%%%%%%%%%%%%%%%%%%%%%%%%%%
\subsection{\texorpdfstring{$t$}{t}-structures}
\label{sec:prelim_t-structures_t_structures}
%%%%%%%%%%%%%%%%%%%%%%%%%%%%%%%%%%%%%%%%%%%%  

A pair of strictly full subcategories $(\mathcal{T}^{\leq 0}, \mathcal{T}^{\geq 0})$ of $\mathcal{T}$ is a \textbf{$t$-structure} if:
\begin{itemize}
    \item $\operatorname{Hom}(A,B) = 0$ for all $A \in \mathcal{T}^{\leq 0}$ and $B \in \mathcal{T}^{\geq 0}[-1]$,
    \item $\mathcal{T}^{\leq 0}[1] \subseteq \mathcal{T}^{\leq 0}$ and $\mathcal{T}^{\geq 0}[-1] \subseteq \mathcal{T}^{\geq 0}$,
    \item for every $E \in \mathcal{T}$, there is a distinguished triangle
    \begin{displaymath}
        \tau^{\leq 0} E \to E \to \tau^{\geq 1} E \to (\tau^{\leq 0} E)[1]
    \end{displaymath}
    with $\tau^{\leq 0} E \in \mathcal{T}^{\leq 0}$ and $\tau^{\geq 1} E \in \mathcal{T}^{\geq 0}[-1]$.
\end{itemize}
This notion was introduced in \cite{Beilinson/Berstein/Deligne/Gabber:2018}. 
The above triangle is called the \textbf{truncation triangle} of $E$ with respect to the $t$-structure. 
For any $n \in \mathbb{Z}$,$(\mathcal{T}^{\leq n}, \mathcal{T}^{\geq n})$ is also a $t$-structure on $\mathcal{T}$ where $\mathcal{T}^{\leq n} := \mathcal{T}^{\leq 0}[-n]$ and $\mathcal{T}^{\geq n} := \mathcal{T}^{\geq 0}[-n]$. 

\begin{example}
    Let $X$ be a quasi-compact quasi-separated algebraic space. 
    Consider the following categories
    \begin{itemize}
        \item $D^{\leq 0}(X)$ consists of objects $E\in D(X)$ such that $\mathcal{H}^j (E)=0$ for $j >0$
        \item $D^{\geq 0}(X)$ consists of objects $E\in D(X)$ such that $\mathcal{H}^j (E)=0$ for $j <0$. 
    \end{itemize}
    We call $(D^{\leq 0}(X), D^{\geq 0}(X))$ the \textbf{standard $t$-structure} on $D(X)$. 
    Since the natural functor $\operatorname{Qcoh}(X) \to \operatorname{Mod}(X)$ is fully faithful and exact \cite[\href{https://stacks.math.columbia.edu/tag/03M1}{Tag 03M1}]{StacksProject}, it follows that the standard $t$-structure induces a $t$-structure on $D_{\operatorname{qc}}(X)$, denoted by $(D^{\leq 0}_{\operatorname{qc}}(X), D^{\geq 0}_{\operatorname{qc}}(X))$.
\end{example}

%%%%%%%%%%%%%%%%%%%%%%%%%%%%%%%%%%%%%%%%%%%%
\subsection{(Pre)aisles}
\label{sec:prelim_t-structures_aisles}
%%%%%%%%%%%%%%%%%%%%%%%%%%%%%%%%%%%%%%%%%%%%  

A strictly full subcategory $\mathcal{A} \subseteq \mathcal{T}$ is an \textbf{aisle} if it is closed under positive shifts and extensions, and the inclusion $\mathcal{A} \hookrightarrow \mathcal{T}$ admits a right adjoint. 
Given a $t$-structure $(\mathcal{T}^{\leq 0}, \mathcal{T}^{\geq 0})$, the subcategory $\mathcal{T}^{\leq 0}$ is an aisle. 
By \cite[\S 1, Proposition]{Keller/Vossieck:1988}, any aisle $\mathcal{A}$ determines a $t$-structure $(\mathcal{A}, \mathcal{A}^\perp[1])$, where
\begin{displaymath}
    \mathcal{A}^\perp := \{ T \in \mathcal{T} \mid \operatorname{Hom}(A,T) = 0 \text{ for all } A \in \mathcal{A} \}.
\end{displaymath}
The subcategories $\mathcal{T}^{\leq 0}$ and $\mathcal{T}^{\geq 0}$ are called the \textbf{aisle} and \textbf{coaisle} of the $t$-structure, respectively. 
More generally, a \textbf{preaisle} is a strictly full subcategory closed under positive shifts and extensions. 
A preaisle $\mathcal{A} \subseteq \mathcal{T}$ is called \textbf{cocomplete} if it is closed under all coproducts in $\mathcal{T}$. 
Given $\mathcal{S}\subseteq \mathcal{T}$, $\overline{\langle \mathcal{S} \rangle}^{(-\infty, 0]}$ is defined to be the smallest cocomplete preaisle containing $\mathcal{S}$. 
By \cite[Theorem 2.3]{Neeman:2021}, if $\mathcal{T}$ is well generated (e.g.\ compactly generated), $\overline{\langle \mathcal{S} \rangle}^{(-\infty, 0]}$ is an aisle whenever $\mathcal{S}$ is essentially small. 
An aisle $U$ on $\mathcal{T}$ is \textbf{compactly generated} when there exists a collection of compact objects $\mathcal{P} \subseteq \mathcal{T}^c$ satisfying $\overline{\langle \mathcal{P} \rangle}^{(-\infty, 0]} = U$. 
We say a $t$-structure is \textbf{compactly generated} if its aisle is as such.  

%%%%%%%%%%%%%%%%%%%%%%%%%%%%%%%%%%%%%%%%%%%%
\subsection{\texorpdfstring{$t$}{t}-exactness}
\label{sec:prelim_t-structures_t_exactness}
%%%%%%%%%%%%%%%%%%%%%%%%%%%%%%%%%%%%%%%%%%%%  

Let $F \colon \mathcal{T} \to \mathcal{T}^\prime$ be an exact functor between triangulated categories equipped with $t$-structures $(\mathcal{T}^{\leq 0}, \mathcal{T}^{\geq 0})$ and $((\mathcal{T}^\prime)^{\leq 0}, (\mathcal{T}^\prime)^{\geq 0})$. 
We say that $F$ is \textbf{right $t$-exact} if $
F(\mathcal{T}^{\leq 0}) \subseteq (\mathcal{T}^\prime)^{\leq 0}$,
and \textbf{left $t$-exact} if 
$F(\mathcal{T}^{\geq 0}) \subseteq (\mathcal{T}^\prime)^{\geq 0}$. 
If both conditions hold, then $F$ is \textbf{$t$-exact}.
At times, we might write a (resp.\ right, left) $t$-exact functor as $F\colon (\mathcal{T},\mathcal{T}^{\leq 0}) \to (\mathcal{T}^\prime ,(\mathcal{T}^\prime)^{\leq 0})$ where only the aisle is made clear.

%%%%%%%%%%%%%%%%%%%%%%%%%%%%%%%%%%%%%%%%%%%%
\subsection{Tensor variants}
\label{sec:prelim_t-structures_tensor}
%%%%%%%%%%%%%%%%%%%%%%%%%%%%%%%%%%%%%%%%%%%%  

We discuss a notion of aisles via tensor actions of triangulated categories. 
See \cite[\S 3]{Stevenson:2013} for the definition of tensor action. 
The following appears in \cite[\S 4]{Hrbek/Lank/LeGros/Pavon:2026}. 
Let $(\mathcal{T},\otimes,\mathbf{1})$ be a rigidly compactly generated tensor triangulated category and $\mathcal{K}$ be a compactly generated triangulated category. 
Suppose $\odot\colon \mathcal{T}\times\mathcal{K}\to \mathcal{K}$ is an action of $\mathcal{T}$ on $\mathcal{K}$. 
Note that $\odot$ is exact and coproduct preserving in both variables. 
By \cite[Lemma 4.6]{Stevenson:2013}, the action restricts to an action $\odot\colon \mathcal{T}^c\times\mathcal{K}^c\to \mathcal{K}^c$.  
Fix a preaisle $\mathcal{P}^{\leq 0}\subseteq\mathcal{T}^c$ satisfying
$\mathbf{1}\in\mathcal{P}^{\leq 0}$ and $\mathcal{P}^{\leq 0}\otimes\mathcal{P}^{\leq 0}\subseteq\mathcal{P}^{\leq 0}$. 
Set $\mathcal{T}^{\leq 0}\colonequals \overline{\langle \mathcal{P}^{\leq 0} \rangle}^{(-\infty, 0]}$ to be the aisle generated by $\mathcal{P}^{\leq 0}$.  
We say that a preaisle $\mathcal{U}\subseteq\mathcal{K}$ is a \textbf{$\odot$-preaisle} if $\mathcal{P}^{\leq 0}\odot \mathcal{U}\subseteq \mathcal{U}$. 
By \cite[Lemma 2.6]{Hrbek/Lank/LeGros/Pavon:2026}, for a cocomplete preaisle  $\mathcal{A}\subseteq\mathcal{K}$, $\mathcal{A}$ is a $\odot$-preaisle if, and only if, $\mathcal{T}^{\leq 0}\odot \mathcal{A}\subseteq\mathcal{A}$. 
Let $\mathcal{S}\subseteq\mathcal{K}$ be a set of objects. 
Denote by $\overline{\langle \mathcal{S} \rangle}^{(-\infty, 0]}_\odot$ the smallest cocomplete $\odot$-preaisle of $\mathcal{K}$ containing $\mathcal{S}$ (which is obtained by intersecting all of them). Given a skeletally small subcategory $\mathcal{S}\subseteq\mathcal{K}$, $\overline{\langle \mathcal{S} \rangle}^{(-\infty, 0]}_\odot = \overline{\langle \mathcal{P}^{\leq 0}\odot\mathcal{S} \rangle}^{(-\infty, 0]}$. 
In fact, $\overline{\langle \mathcal{S} \rangle}^{(-\infty, 0]}_\odot$ is an aisle. 
See \cite[Lemma 2.8]{Hrbek/Lank/LeGros/Pavon:2026}.

%%NOTE: If $X$ is an affine Noetherian scheme, then every aisle on $D_{\operatorname{qc}}(X)$ is an $\otimes$-aisle \cite[Proposition 3.8]{Dubey/Sahoo:2023}. This fact for the affine scheme case comes from the following lemma coupled with the discussion at start of \cite[\S 2.3]{Canonaco/Neeman/Stellari:2025}. Loc.\ cit.\ tells us that the aisle for the standard $t$-structure on $D_{\operatorname{qc}}(X)$ for an affine scheme $X$ is compactly generated by a single object.

\begin{lemma}
    \label{lem:dubey_sahoo_corrected}
    Let $(\mathcal{T},\otimes,\mathbf{1})$ be a rigidly compactly generated tensor triangulated category.  Suppose $\mathcal{T}^{\leq 0}$ is a compactly generated preaisle. Assume that $\mathbf{1} \in \mathcal{P}^{\leq 0}:= \mathcal{T}^c \cap \mathcal{T}^{\leq 0}$. For any $\mathcal{S}\subseteq \mathcal{T}^c$ and $A\in \overline{\langle \mathcal{S} \rangle}^{(-\infty,0]}_{\otimes}$, there is a distinguished triangle 
    \begin{displaymath}
        \bigsqcup_{i\geq 0} A_i \to A \to (\bigsqcup_{i\geq 0} A_i)[1] \to (\bigsqcup_{i\geq 0} A_i)[1]
    \end{displaymath}
    where $A_i$ are $i$-fold extensions of small coproducts of nonnegative shifts of objects in $\mathcal{P}^{\leq 0}\otimes \mathcal{S}$.
\end{lemma}

\begin{proof}
    This is exactly \cite[Proposition 3.11]{Dubey/Sahoo:2023} with the correction that $\mathbf{1} \in \mathcal{P}^{\leq 0}:= \mathcal{T}^c \cap \mathcal{T}^{\leq 0}$ is added as a hypothesis.
\end{proof}

%%%%%%%%%%%%%%%%%%%%%%%%%%%%%%%%%%%%%%%%%%%%
\subsection{Thomason filtrations}
\label{sec:prelim_Thomason}
%%%%%%%%%%%%%%%%%%%%%%%%%%%%%%%%%%%%%%%%%%%

We recall Thomason filtrations on a topological space $X$.  
Recall that $Z\subseteq X$ is called a \textbf{Thomason subset} of $X$ if $Z=\cup_{\alpha \in \Lambda} Z_{\alpha}$ where each $Z_{\alpha}$ is closed in $X$ and $X\setminus Z_{\alpha}$ is quasi-compact. 
Denote by $\operatorname{Thom}(X)$ the collection of Thomason subsets of $X$. 
If $X$ is a Noetherian topological space, then specialization closed subsets coincide with Thomason subsets. 
A function $\phi \colon \mathbb{Z} \to \operatorname{Thom}(X)$ is called a \textbf{Thomason filtration} on $X$ if $\phi(n+1)\subseteq \phi (n)$ for each $n\in \mathbb{Z}$. 
If $X=|X^\prime|$ where $X^\prime$ is an algebraic space, then we abuse language and call these `Thomason filtrations on $X$'.

%%%%%%%%%%%%%%%%%%%%%%%%%%%%%%%%%%%%%
\section{Pseudoapproximation}
\label{sec:pseudoapproximation}
%%%%%%%%%%%%%%%%%%%%%%%%%%%%%%%%%%%%%

We introduce a weak version of \cite[\href{https://stacks.math.columbia.edu/tag/08HH}{Tag 08HH}]{StacksProject}.
In particular, \Cref{lem:represent_bounded_above_pseudocoherent,lem:perf_up_to_summand,lem:lift_roofs,lem:lift_up_to_bundle,lem:perf_up_to_summand_support,lem:almost_lift_psuedoapprox} are cohomological sensitive variations of \cite[\href{https://stacks.math.columbia.edu/tag/08EC}{Tag 08EC}]{StacksProject}.

\begin{definition}
    \label{def:approximation}
    Let $X$ be an algebraic space. 
    Consider a pair $(Z, E)$  consisting of a closed subset $Z\subseteq|X|$ and a complex $E\in D^{\leq 0}_{\operatorname{qc},Z}(X)$ such that $\mathcal{H}^0 (E)$ is of finite type. 
    \begin{enumerate}
        \item  We say that \textbf{pseudoapproximation by perfect complexes holds for $(Z, E)$} if there exist $P\in D^{\leq 0}_{\operatorname{qc},Z}(X) \cap \operatorname{Perf}(X)$ and 
        $P\to E$ such that $\mathcal{H}^0 (P) \to \mathcal{H}^0 (E)$ is surjective. 
        \item We say that \textbf{pseudoapproximation by perfect complexes holds
        for $X$} if it holds for any pair $(Z, E)$ where $Z\subseteq |X|$ is closed with quasi-compact complement and $E\in D^{\leq 0}_{\operatorname{qc},Z}(X)$ has $\mathcal{H}^0 (E)$ of finite type.
    \end{enumerate}
\end{definition}

\begin{proposition}
    \label{prop:approx_by_comorphismcts_implies_Thomason}
    Let $X$ be a quasi-compact quasi-separated algebraic space. 
    If $X$ satisfies pseudoapproximation by perfect complexes, then $D^{\leq 0}_{\operatorname{qc},Z}(X)$ is compactly generated for any closed $Z\subseteq |X|$ with quasi-compact complement.
\end{proposition}

\begin{proof}
    Fix a closed subset $Z$ of $|X|$ with quasi-compact complement.
    By construction, 
    \begin{displaymath}
        \overline{\langle \operatorname{Perf}_Z (X) \cap D^{\leq 0}_{\operatorname{qc},Z}(X) \rangle}^{(-\infty,0]} \subseteq D^{\leq 0}_{\operatorname{qc},Z}(X).
    \end{displaymath} 
    We prove the reverse inclusion. 
    This is equivalent to showing that
    \begin{displaymath}
        (\overline{\langle \operatorname{Perf}_Z (X) \cap D^{\leq 0}_{\operatorname{qc},Z}(X) \rangle}^{(-\infty,0]})^\perp \subseteq D^{\geq 1}_{\operatorname{qc},Z}(X),
    \end{displaymath}
    which will be done by contradiction.

    Choose a compact generator $G$ for $D_{\operatorname{qc},Z}(X)$, which exists by \cite[\href{https://stacks.math.columbia.edu/tag/0AEC}{Tag 0AEC}]{StacksProject}. 
    If needed, we can shift so that $G\in D^{\leq 0}_{\operatorname{qc},Z}(X)$. 
    Before proving the claim, we make the observation that
    \begin{displaymath}
        (\overline{\langle \operatorname{Perf}_Z (X) \cap D^{\leq 0}_{\operatorname{qc},Z}(X) \rangle}^{(-\infty,0]})^\perp \subseteq D^+_{\operatorname{qc},Z}(X).
    \end{displaymath} 
    Indeed, if $G\in D^{\leq 0}_{\operatorname{qc},Z}(X)$, then
    \begin{displaymath}
        \overline{\langle G \rangle}^{(-\infty,0]} \subseteq \overline{\langle \operatorname{Perf}_Z (X) \cap D^{\leq 0}_{\operatorname{qc},Z}(X) \rangle}^{(-\infty,0]}.
    \end{displaymath} 
    By \cite[Lemma A.8]{GuisadoVillaalgordo/Lank:2026}, there exists $n\geq 0$ such that 
    \begin{displaymath}
        D^{\leq -n}_{\operatorname{qc},Z}(X) 
        \subseteq \overline{\langle G \rangle}^{(-\infty,0]} 
        \subseteq D^{\leq n}_{\operatorname{qc},Z}(X).
    \end{displaymath}
    Consequently, 
    \begin{displaymath}
        D^{\leq -n}_{\operatorname{qc},Z}(X) \subseteq \overline{\langle \operatorname{Perf}_Z (X) \cap D^{\leq 0}_{\operatorname{qc},Z}(X) \rangle}^{(-\infty,0]},
    \end{displaymath}
    and so,
    \begin{displaymath}
        (\overline{\langle \operatorname{Perf}_Z (X) \cap D^{\leq 0}_{\operatorname{qc},Z}(X) \rangle}^{(-\infty,0]})^\perp 
        \subseteq D^{\geq -n+1}_{\operatorname{qc},Z}(X) 
        \subseteq D^+_{\operatorname{qc},Z}(X).
    \end{displaymath}

    Now, we prove that $D^{\leq 0}_{\operatorname{qc},Z}(X)$ is compactly generated. 
    By assumption, 
    \begin{displaymath}
        (\overline{\langle \operatorname{Perf}_Z (X) \cap D^{\leq 0}_{\operatorname{qc},Z}(X) \rangle}^{(-\infty,0]})^\perp \not\subseteq D^{\geq 1}_{\operatorname{qc},Z}(X).
    \end{displaymath} 
    This gives us some nonzero 
    \begin{displaymath}
        E\in (\overline{\langle \operatorname{Perf}_Z (X) \cap D^{\leq 0}_{\operatorname{qc},Z}(X) \rangle}^{(-\infty,0]})^\perp
    \end{displaymath}
    such that $E \not\in D^{\geq 1}_{\operatorname{qc},Z}(X)$. 
    As 
    \begin{displaymath}
        (\overline{\langle \operatorname{Perf}_Z (X) \cap D^{\leq 0}_{\operatorname{qc},Z}(X) \rangle}^{(-\infty,0]})^\perp \subseteq D^+_{\operatorname{qc},Z}(X)
    \end{displaymath}
    there exists a minimal $s\leq 0$ such that $\mathcal{H}^s (E)\not=0$. 
    By \cite[\href{https://stacks.math.columbia.edu/tag/0829}{Tag 0829}]{StacksProject}, $\mathcal{H}^s (E)$ is a directed colimit of its finite type $\mathcal{O}_{X}$-submodules $E_i$. 
    Since $\mathcal{H}^s (E)$ is supported on $Z$, so is each $E_i$ (see \cite[\href{https://stacks.math.columbia.edu/tag/056J}{Tags 056J} \& \href{https://stacks.math.columbia.edu/tag/07TY}{07TY}]{StacksProject}). 
    We can find some $E_t$ such that $\operatorname{Hom}(E_t,\mathcal{H}^s (E))\not=0$. 
    
    As $X$ satisfies pseudoapproximation by perfect complexes, there exists 
    \begin{displaymath}
        P\in D^{\leq 0}_{\operatorname{qc},Z}(X) \cap \operatorname{Perf}_Z (X)
    \end{displaymath}
    and $P\to E_t$ such that $\mathcal{H}^0 (P) \to \mathcal{H}^0 (E_t)$ is surjective. 
    Since $P[-s]\in D^{\leq s}_{\operatorname{qc},Z}(X)$ and $E\in D^{\geq s}_{\operatorname{qc},Z}(X)$, we obtain from properties of truncations, 
    \begin{displaymath}
        \begin{aligned}
            \operatorname{Hom}(P[-s],E) 
            &\cong \operatorname{Hom}(\mathcal{H}^s (P[-s]), \mathcal{H}^s (E))
            \\&\cong \operatorname{Hom}(\mathcal{H}^0 (P), \mathcal{H}^s (E)) 
            \\&\not\cong 0.
        \end{aligned}
    \end{displaymath}
    Then $\mathcal{H}^s (P[-s]) = \mathcal{H}^0 (P)$ and the composite
    \begin{displaymath}
        \mathcal{H}^0 (P) \to E_t \to \mathcal{H}^s (E)
    \end{displaymath}
    is nonzero. 
    Thus, we obtain a contradiction.
    % Indeed, the morphisms provided above imply the composition 
    % \begin{displaymath}
    %     \mathcal{H}^0 (P) \to E_t \to \mathcal{H}^s (E)
    % \end{displaymath}
    % is nonzero. 
    % Since 
    % \begin{displaymath}
    %     P[-s]\in \overline{\langle \operatorname{Perf}_Z (X) \cap D^{\leq 0}_{\operatorname{qc},Z}(X) \rangle}^{(-\infty,0]}
    % \end{displaymath}
    % and 
    % \begin{displaymath}
    %     E\in (\overline{\langle \operatorname{Perf}_Z (X) \cap D^{\leq 0}_{\operatorname{qc},Z}(X) \rangle}^{(-\infty,0]})^\perp,
    % \end{displaymath}
    % we obtain a contradiction.
\end{proof}

\begin{lemma}
    \label{lem:quasi-affine_pseudoapprox}
    Let $X=\operatorname{Spec}(R)$ be an affine scheme. 
    Then $X$ satisfies pseudoapproximation by perfect complexes.
\end{lemma}

\begin{proof}
    Applying \cite[\href{https://stacks.math.columbia.edu/tag/071Q}{Tag 071Q}]{StacksProject}, we can transfer the problem to the small Zariski site of $X$.
    Fix a closed subset $Z\subseteq |X|$ with quasi-compact complement. 
    Choose $E\in D^{\leq 0}_{\operatorname{qc},Z}(X)$ such that $\mathcal{H}^0 (E)$ is of finite type. 
    Let $j\colon U \to X$ be the open immersion with $|U|:= |X|\setminus Z$. 
    Since $E\in D^{\leq 0}_{\operatorname{qc},Z}(X)$, one has $\mathbf{L}j^\ast E\cong 0$. 

    Given that $U$ is quasi-compact, there exist $f_1,\ldots,f_r \in R$ such that $U=\cup^r_{i=1} D(f_i)$.
    Because $X$ is affine and $\mathcal{H}^0 (E)$ is of finite type, we can choose global sections $s_1,\ldots,s_c \in H^0 (X,\mathcal{H}^0 (E))$ which generate $H^0 (X,\mathcal{H}^0 (E))$ as an $R$-module. 
    By \cite[\href{https://stacks.math.columbia.edu/tag/06Z0}{Tag 06Z0}]{StacksProject}, we may regard these as elements of $H^0(X,E)$. 
    For each $a=1,\ldots,c$, \cite[\href{https://stacks.math.columbia.edu/tag/08E3}{Tag 08E3}]{StacksProject} gives an integer $N_a \geq 0$ and a morphism $g_a \colon K_{N_a} \to E$ such that $s_a$ is in the image of the $H^0(X,K_{N_a}) \to H^0(X,E)$ where $K_{N_a}$ is the object of $D(X)$ corresponding to the Koszul complex $K^\bullet (f_1^{N_a},\ldots,f_r^{N_a})$ via the equivalence \cite[\href{https://stacks.math.columbia.edu/tag/06Z0}{Tag 06Z0}]{StacksProject}. 
    The hypothesis of \cite[\href{https://stacks.math.columbia.edu/tag/08E3}{Tag 08E3}]{StacksProject} holds because $\mathbf{L}j^\ast E\cong 0$. 
    Set $P:= \oplus^c_{a=1} K_{N_a}$. 
    Consider the natural morphism $g\colon P\to E$ obtained from the $g_a$.
    %%NOTE: Use def of coproducts and universal properties

    Each $K_{N_a}$ is a bounded complex of finite locally free sheaves whose cohomology is concentrated in nonpositive degrees, and its cohomology is supported on $Z$. 
    The image of $H^0(X,P)\to H^0(X,E)$ contains all of the $s_a$, and hence is surjective. 
    Since $X$ is affine, \cite[\href{https://stacks.math.columbia.edu/tag/06Z0}{Tag 06Z0}]{StacksProject} identifies these groups with the global sections of the corresponding degree zero cohomology sheaves. 
    Hence, the induced morphism $\mathcal{H}^0 (P) \to \mathcal{H}^0 (E)$ is surjective. 
    This proves that pseudoapproximation holds for $(Z,E)$, and hence completes the proof.
\end{proof}

\begin{lemma}
    \label{lem:represent_bounded_above_pseudocoherent}
    Let $X$ be an affine scheme, $j\colon U \to X$ a quasi-compact open immersion, $N\in \mathbb{Z}$, and $E\in D^{\leq N}(U)$ be pseudocoherent. 
    There exists $E^\prime\in D^{\leq N}(X)$ whose components are finite free $\mathcal{O}_X$-modules such that $\mathbf{L}j^\ast E^\prime \cong E$. 
\end{lemma}

\begin{proof}
    This follows from the argument of \cite[\href{https://stacks.math.columbia.edu/tag/08EE}{Tag 08EE}]{StacksProject} but we clarify how. 
    Applying \cite[\href{https://stacks.math.columbia.edu/tag/071Q}{Tag 071Q}]{StacksProject}, we can transfer the problem to the small Zariski site of $X$.
    By \cite[\href{https://stacks.math.columbia.edu/tag/08E5}{Tag 08E5}]{StacksProject}, we see that $E$ is an object of $D_{\operatorname{qc}}(\mathcal{O}_ U)$. 
    By \cite[\href{https://stacks.math.columbia.edu/tag/08ED}{Tag 08ED}]{StacksProject}, we may assume $E = \mathbf{L}j^\ast E^\prime$ for some object $E^\prime$ of $D_{\operatorname{qc}}(\mathcal{O}_X)$. Write $X = \operatorname{Spec}(A)$. By \cite[\href{https://stacks.math.columbia.edu/tag/06Z0}{Tag 06Z0}]{StacksProject}, we can find a complex $M^\bullet $ of $A$-modules whose associated complex of $\mathcal{O}_ X$-modules is a representative of $E^\prime$. 

    Choose $f_1, \ldots , f_ r \in A$ such that $U = D(f_1) \cup \ldots \cup D(f_ r)$. 
    By \cite[\href{https://stacks.math.columbia.edu/tag/08E7}{Tag 08E7}]{StacksProject}, the complexes $M^\bullet _{f_ j}$ are pseudo-coherent complexes of $A_{f_ j}$-modules. 
    Let $n$ be an integer. 
    Assume we have a map of complexes $\alpha \colon F^\bullet \to M^\bullet $ where $F^\bullet $ is bounded above, $F^ i = 0$ for $i < n$, each $F^ i$ is a finite free $A$-module, such that 
    \begin{displaymath}
        H^ i(\alpha _{f_ j}) \colon H^ i(F^\bullet _{f_ j}) \to H^ i(M^\bullet _{f_ j}) 
    \end{displaymath}
    is an isomorphism for $i > n$ and surjective for $i = n$. Picture 
    \begin{displaymath}
        % https://q.uiver.app/#q=WzAsNyxbMSwwLCJGXm4iXSxbMiwwLCJGXntuKzF9Il0sWzMsMCwiXFxjZG90cyJdLFsxLDEsIk1ebiJdLFsyLDEsIk1ee24rMX0iXSxbMCwxLCJNXntuLTF9Il0sWzMsMSwiXFxjZG90cy4iXSxbMCwxXSxbMSwyXSxbMCwzLCJcXGFscGhhIiwyXSxbMSw0LCJcXGFscGhhIiwyXSxbMyw0XSxbNSwzXSxbNCw2XV0=
        \begin{tikzcd}
            & {F^n} & {F^{n+1}} & \cdots \\
            {M^{n-1}} & {M^n} & {M^{n+1}} & {\cdots.}
            \arrow[from=1-2, to=1-3]
            \arrow["\alpha"', from=1-2, to=2-2]
            \arrow[from=1-3, to=1-4]
            \arrow["\alpha"', from=1-3, to=2-3]
            \arrow[from=2-1, to=2-2]
            \arrow[from=2-2, to=2-3]
            \arrow[from=2-3, to=2-4]
        \end{tikzcd}
    \end{displaymath}
    Since each $\mathcal{H}^i (M^\bullet _{f_ j})\cong 0$ for $i> N$, we can find such a map for $n$. 
    In fact, we can start the construction in degree $N+1$ (i.e.\ the induction argument regards $n\leq N+1$), which is where we have applied the hypothesis that $E\in D^{\leq N}_{\operatorname{qc}}(U)$. 
    Indeed, we can form the diagram
    \begin{displaymath}
        % https://q.uiver.app/#q=WzAsNyxbMSwwLCJGXntOKzF9Il0sWzIsMCwiRl57bisyfSJdLFszLDAsIlxcY2RvdHMiXSxbMSwxLCJNXm4iXSxbMiwxLCJNXntuKzF9Il0sWzAsMSwiTV57bi0xfSJdLFszLDEsIlxcY2RvdHMiXSxbMCwxXSxbMSwyXSxbMCwzLCIwIiwyXSxbMSw0LCIwIiwyXSxbMyw0XSxbNSwzXSxbNCw2XV0=
        \begin{tikzcd}
            & {F^{N+1}} & {F^{N+2}} & \cdots \\
            {M^{N-1}} & {M^N} & {M^{N+1}} & \cdots
            \arrow[from=1-2, to=1-3]
            \arrow["0"', from=1-2, to=2-2]
            \arrow[from=1-3, to=1-4]
            \arrow["0"', from=1-3, to=2-3]
            \arrow[from=2-1, to=2-2]
            \arrow[from=2-2, to=2-3]
            \arrow[from=2-3, to=2-4]
        \end{tikzcd}
    \end{displaymath}
    where $F^{N+1} \cong F^{N+2} \cong 0$. 
    Since $\mathcal{H}^i (M^\bullet _{f_ j})\cong 0$ for $i> N$, the vertical morphisms induced isomorphisms in cohomological degrees $i>N$. 
    By induction on $n\leq N+1$, we are going to extend this to a map of complexes $F^\bullet \to M^\bullet $ such that $H^ i(\alpha _{f_ j})$ is an isomorphism for all $i$. 
    The desired claim will follow by taking $\widetilde{F^\bullet }$. 

    The induction step will be to extend the diagram above by adding $F^{n - 1}$. 
    Let $C^\bullet $ be the cone on $\alpha $ \cite[\href{https://stacks.math.columbia.edu/tag/014E}{Tag 014E}]{StacksProject}. 
    The long exact sequence of cohomology shows that $H^ i(C^\bullet _{f_ j}) = 0$ for $i \geq n$. 
    By \cite[\href{https://stacks.math.columbia.edu/tag/064R}{Tag 064R}]{StacksProject}, we see that $C^\bullet _{f_ j}$ is $(n - 1)$-pseudo-coherent. 
    By \cite[\href{https://stacks.math.columbia.edu/tag/064S}{Tag 064S}]{StacksProject}, we see that $H^{n - 1}(C^\bullet _{f_ j})$ is a finite $A_{f_ j}$-module. 
    Choose a finite free $A$-module $F^{n - 1}$ and an $A$-module morphism $\beta \colon F^{n - 1} \to C^{n - 1}$ such that the composition $F^{n - 1} \to C^{n - 1} \to C^ n$ is zero and such that $F^{n - 1}_{f_ j}$ surjects onto $H^{n - 1}(C^\bullet _{f_ j})$ (e.g.\ clear denominators). 
    Since $C^{n - 1} = M^{n - 1} \oplus F^ n$ we can write $\beta = (\alpha ^{n - 1}, -d^{n - 1})$. 
    The vanishing of the composition $F^{n - 1} \to C^{n - 1} \to C^ n$ implies these maps fit into a morphism of complexes 
    \begin{displaymath}
        % https://q.uiver.app/#q=WzAsOSxbMiwwLCJGXm4iXSxbMywwLCJGXntuKzF9Il0sWzQsMCwiXFxjZG90cyJdLFsyLDEsIk1ebiJdLFszLDEsIk1ee24rMX0iXSxbMSwxLCJNXntuLTF9Il0sWzQsMSwiXFxjZG90cy4iXSxbMSwwLCJGXntuLTF9Il0sWzAsMSwiXFxjZG90cyJdLFswLDFdLFsxLDJdLFswLDMsIlxcYWxwaGEiLDJdLFsxLDQsIlxcYWxwaGEiLDJdLFszLDRdLFs1LDNdLFs0LDZdLFs3LDAsImRee24tMX0iXSxbNyw1LCJcXGFscGhhXntuLTF9IiwyXSxbOCw1XV0=
        \begin{tikzcd}
            & {F^{n-1}} & {F^n} & {F^{n+1}} & \cdots \\
            \cdots & {M^{n-1}} & {M^n} & {M^{n+1}} & {\cdots.}
            \arrow["{d^{n-1}}", from=1-2, to=1-3]
            \arrow["{\alpha^{n-1}}"', from=1-2, to=2-2]
            \arrow[from=1-3, to=1-4]
            \arrow["\alpha"', from=1-3, to=2-3]
            \arrow[from=1-4, to=1-5]
            \arrow["\alpha"', from=1-4, to=2-4]
            \arrow[from=2-1, to=2-2]
            \arrow[from=2-2, to=2-3]
            \arrow[from=2-3, to=2-4]
            \arrow[from=2-4, to=2-5]
        \end{tikzcd}
    \end{displaymath}
    Moreover, these maps define a morphism of distinguished triangles
    \begin{displaymath}
        % https://q.uiver.app/#q=WzAsOCxbMSwwLCIoRl57bi0xfSBcXHRvIFxcY2RvdHMpIl0sWzIsMCwiRl57bi0xfSJdLFsxLDEsIk1eXFxidWxsZXQiXSxbMiwxLCJDXlxcYnVsbGV0Il0sWzAsMSwiKEZebiBcXHRvIFxcY2RvdHMpIl0sWzAsMCwiKEZebiBcXHRvIFxcY2RvdHMpIl0sWzMsMCwiKEZebiBcXHRvIFxcY2RvdHMpWzFdIl0sWzMsMSwiKEZebiBcXHRvIFxcY2RvdHMpWzFdLiJdLFswLDFdLFswLDJdLFsxLDMsIlxcYmV0YSIsMl0sWzIsM10sWzQsMl0sWzUsMF0sWzUsNF0sWzEsNl0sWzMsN10sWzYsN11d
        \begin{tikzcd}
            {(F^n \to \cdots)} & {(F^{n-1} \to \cdots)} & {F^{n-1}[1-n]} & {(F^n \to \cdots)[1]} \\
            {(F^n \to \cdots)} & {M^\bullet} & {C^\bullet} & {(F^n \to \cdots)[1].}
            \arrow[from=1-1, to=1-2]
            \arrow[from=1-1, to=2-1]
            \arrow[from=1-2, to=1-3]
            \arrow[from=1-2, to=2-2]
            \arrow[from=1-3, to=1-4]
            \arrow["\beta"', from=1-3, to=2-3]
            \arrow[from=1-4, to=2-4]
            \arrow[from=2-1, to=2-2]
            \arrow[from=2-2, to=2-3]
            \arrow[from=2-3, to=2-4]
        \end{tikzcd}
    \end{displaymath}
    Hence our choice of $\beta $ implies that the map of complexes $(F^{n-1} \to \ldots ) \to M^\bullet $ induces an isomorphism on cohomology localized at $f_ j$ in degrees $\geq n$ and a surjection in degree $n - 1$. 
    This finishes the proof.
\end{proof}

\begin{lemma}
    \label{lem:lift_up_to_bundle}
    Let $X=\operatorname{Spec}(R)$ be an affine scheme, $j\colon U \to X$ be a quasi-compact open immersion, and $P\in D^{\leq 0}_{\operatorname{qc}}(U) \cap \operatorname{Perf}(U)$. 
    There exist a finite locally free sheaf $G$ on $U$, $P^\prime \in \operatorname{Perf}(X) \cap D^{\leq 0}_{\operatorname{qc}}(X)$, and $s\geq 0$ such that $\mathbf{L}j^\ast P^\prime \cong P\oplus G[s]$. 
\end{lemma}

\begin{proof}
    We modify the proof of \cite[\href{https://stacks.math.columbia.edu/tag/08EG}{Tag 08EG}]{StacksProject}.
    Applying \cite[\href{https://stacks.math.columbia.edu/tag/071Q}{Tag 071Q}]{StacksProject}, we can transfer the problem to the small Zariski site of $X$.
    Recall that a perfect complex is pseudocoherent \cite[\href{https://stacks.math.columbia.edu/tag/08CQ}{Tag 08CQ}]{StacksProject}. 
    By \Cref{lem:represent_bounded_above_pseudocoherent}, there exists $E^\bullet \in D^{\leq 0}(R)$ of finite free $R$-modules such that $P\cong \mathbf{L}j^\ast E^\bullet$. 
    Furthermore, from \cite[\href{https://stacks.math.columbia.edu/tag/08CQ}{Tag 08CQ}]{StacksProject}, $U$ being quasi-compact implies $P$ has tor amplitude $[a,b]$ for some $a,b\in \mathbb{Z}$. 
    Choose $r< a$. 
    Set 
    \begin{displaymath}
        G^\prime:= \ker (E^r \xrightarrow{d^r} E^{r+1}) = \operatorname{im}(E^{r-1} \xrightarrow{d^{r-1}} E^r).
    \end{displaymath}
    As $P$ has finite tor amplitude in $[a,b]$, we see that $G := \mathbf{L}j^\ast G^\prime$ is flat and of finite presentation.
    Hence, $G$ is finite locally free \cite[\href{https://stacks.math.columbia.edu/tag/05P2}{Tag 05P2}]{StacksProject}.
    Observe that 
    \begin{displaymath}
        G \to j^\ast E^r \xrightarrow{j^\ast d^r} j^\ast E^{r+1} \xrightarrow{j^\ast d^{r+1}} \cdots 
    \end{displaymath}
    is a strictly perfect complex on $U$ which represents $P$. 
    Furthermore, the complex 
    \begin{displaymath}
        Q := E^r \xrightarrow{d^r} E^{r+1} \xrightarrow{d^{r+1}} \cdots
    \end{displaymath}
    is a perfect complex on $X$. 
    As $E^\bullet \in D^{\leq 0}(R)$, it follows that $Q\in D^{\leq 0}(R)$. 
    Taking stupid truncations yields a distinguished triangle
    \begin{displaymath}
        \mathbf{L}j^\ast Q \to P \to G [-r -1] \to \mathbf{L}j^\ast Q[1].
    \end{displaymath}
    By \cite[\href{https://stacks.math.columbia.edu/tag/09M4}{Tag 09M4}]{StacksProject}, there exists $r \ll 0$ such that $P \to G [-r -1]$ is zero, and so $\mathbf{L}j^\ast Q \cong G[-r-2]\oplus P$ (after decreasing $r$, if necessary). 
    This completes the proof.
\end{proof}

\begin{lemma}
    \label{lem:lift_roofs}
    Let $X$ be an affine scheme, $j\colon U \to X$ be a quasi-compact open immersion, and $E,E^\prime \in D^{\leq 0}_{\operatorname{qc}}(X)$ with $E$ a perfect complex supported on a closed subset $Z \subseteq |X|$. 
    For every $\alpha \colon \mathbf{L}j^\ast E \to \mathbf{L}j^\ast E^\prime$, there exist $E_1 \in \operatorname{Perf}(X) \cap D^{\leq 0}_{\operatorname{qc},Z}(X)$ and morphisms
    \begin{displaymath}
        % https://q.uiver.app/#q=WzAsMyxbMSwwLCJFXzEiXSxbMCwxLCJFIl0sWzIsMSwiRV5cXHByaW1lIl0sWzAsMSwiXFxiZXRhIiwyXSxbMCwyLCJcXGdhbW1hIl1d
        \begin{tikzcd}
            & {E_1} & \\
            E && {E^\prime}
            \arrow["\beta"', from=1-2, to=2-1]
            \arrow["\gamma", from=1-2, to=2-3]
        \end{tikzcd}
    \end{displaymath}
    such that $\mathbf{L}j^\ast \beta$ is an isomorphism and $\alpha = \mathbf{L}j^\ast \gamma \circ (\mathbf{L}j^\ast \beta)^{-1}$.
\end{lemma}

\begin{proof}
    This comes from the argument of \cite[\href{https://stacks.math.columbia.edu/tag/08EH}{Tag 08EH}]{StacksProject}. 
    Apply \cite[\href{https://stacks.math.columbia.edu/tag/071Q}{Tag 071Q}]{StacksProject} to transfer the problem to the small Zariski site.
    In particular, loc.\ cit.\ says we can take $E_1 := E \otimes^{\mathbf{L}} I^\bullet(f_1^e,\ldots,f^e_r)$ for some $e\geq 0$ where $U= \cup^r_{i=1} D(f_i)$. 
    Here, $I:=I^\bullet(f_1^e,\ldots,f^e_r)$ is the nonpositively graded perfect complex constructed in the proof of \cite[\href{https://stacks.math.columbia.edu/tag/08EH}{Tag 08EH}]{StacksProject}, see \cite[\href{https://stacks.math.columbia.edu/tag/08CX}{Tag 08CX}]{StacksProject}. 
    Since $E$ and $I^\bullet(f_1^e,\ldots,f^e_r)$ belong to $D^{\leq 0}_{\operatorname{qc}}(X)$, it follows that $ E \otimes^{\mathbf{L}} I^\bullet(f_1^e,\ldots,f^e_r) \in D^{\leq 0}_{\operatorname{qc}}(X)$. 
    Moreover, if $E$ is supported on $Z$, then so is $E \otimes^{\mathbf{L}} I^\bullet(f_1^e,\ldots,f^e_r)$. 
    Thus, the desired claim follows. 
\end{proof}

\begin{lemma}
    \label{lem:perf_up_to_summand}
    Let $X$ be an affine scheme, $j\colon U \to X$ be a quasi-compact open immersion and $P\in \operatorname{Perf}(U) \cap D^{\leq 0}_{\operatorname{qc}}(U)$. 
    There exists $P^\prime \in \operatorname{Perf}(X) \cap D^{\leq 0}_{\operatorname{qc}}(X)$ such that $\mathbf{L}j^\ast P^\prime \cong P \oplus P[1]$.
\end{lemma}

\begin{proof}
    By \Cref{lem:lift_up_to_bundle}, there exist a finite locally free sheaf $G$ on $U$, $P^\prime \in D^{\leq 0}_{\operatorname{qc}}(X)$, and $s\geq 0$ such that $\mathbf{L}j^\ast P^\prime \cong P\oplus G[s]$.
    Moreover, from \Cref{lem:lift_roofs}, there exist $P_1\in \operatorname{Perf}(X)\cap D^{\leq 0}_{\operatorname{qc}}(X)$ and $\alpha \colon P_1 \to P^\prime$ such that $\mathbf{L}j^\ast P_1 \cong \mathbf{L}j^\ast P^\prime$ and $\mathbf{L}j^\ast \alpha$ is the morphism 
    \begin{displaymath}
        \begin{pmatrix}
        1_{G[s]} & 0 \\
        0 & 0 
        \end{pmatrix}
        \colon G[s]\oplus P \to G[s]\oplus P.
    \end{displaymath}
    %%NOTE: Note that we can choose $P_1$ so that $\operatorname{supp}(P_1) = \operatorname{supp}(P^\prime)$ but we don't need that at this step in the proof.
    As $P_1$ and $P^\prime$ are in $\operatorname{Perf}(X)\cap D^{\leq 0}_{\operatorname{qc}}(X)$, it follows that $\operatorname{cone}(\alpha)\in \operatorname{Perf}(X) \cap D^{\leq 0}_{\operatorname{qc}}(X)$. 
    Moreover, $\mathbf{L}j^\ast \operatorname{cone}(\alpha) \cong P \oplus P[1]$, which completes the proof.
\end{proof}

\begin{lemma}
    \label{lem:perf_up_to_summand_support}
    Let $X$ be an affine scheme, $j\colon U \to X$ be a quasi-compact open immersion, $T\subseteq |X|$ be closed with $|X|\setminus T$ quasi-compact, and $F\in \operatorname{Perf}(U) \cap D^{\leq 0}_{\operatorname{qc},T \cap U}(U)$. 
    There exists $E \in \operatorname{Perf}(X) \cap D^{\leq 0}_{\operatorname{qc},T}(X)$ such that $\mathbf{L}j^\ast E \cong F \oplus F[1]$.
\end{lemma}

\begin{proof}
    Applying \cite[\href{https://stacks.math.columbia.edu/tag/071Q}{Tag 071Q}]{StacksProject}, we can transfer the problem to the small Zariski site.
    We modify the proof of \cite[\href{https://stacks.math.columbia.edu/tag/08EK}{Tag 08EK}]{StacksProject}.
    Say $T = V(g_1, \ldots , g_ s)$. 
    After replacing $g_ j$ by a power we may assume multiplication by $g_ j$ is zero on $F$, see \cite[\href{https://stacks.math.columbia.edu/tag/08EJ}{Tag 08EJ}]{StacksProject}. 
    Choose $E\in D^{\leq 0}_{\operatorname{qc}}(X)$ as in \Cref{lem:perf_up_to_summand}.
    Note that $g_ j \colon E \to E$ restricts to zero on $U$. Choose a distinguished triangle of objects in $D^{\leq 0}_{\operatorname{qc}}(X)$,
    \begin{displaymath}
        E \xrightarrow {g_1} E \to C_1 \to E[1].
    \end{displaymath}
    By \cite[\href{https://stacks.math.columbia.edu/tag/05QT}{Tag 05QT}]{StacksProject}, the object $C_1$ restricts to $F \oplus F[1] \oplus F[1] \oplus F[2]$ on $U$. 
    Moreover, $g_1 \colon C_1 \to C_1$ has square zero by \cite[\href{https://stacks.math.columbia.edu/tag/05QP}{Tag 05QP}]{StacksProject}.
    Namely, the diagram 
    \begin{displaymath}
        % https://q.uiver.app/#q=WzAsNixbMCwwLCJFIl0sWzAsMSwiRSJdLFsxLDAsIkNfMSJdLFsxLDEsIkNfMSJdLFsyLDAsIkVbMV0iXSxbMiwxLCJFWzFdIl0sWzAsMSwiMCIsMl0sWzAsMl0sWzEsM10sWzIsMywiZ18xIiwyXSxbMiw0XSxbMyw1XSxbNCw1LCIwIiwyXV0=
        \begin{tikzcd}
            E & {C_1} & {E[1]} \\
            E & {C_1} & {E[1]}
            \arrow[from=1-1, to=1-2]
            \arrow["0"', from=1-1, to=2-1]
            \arrow[from=1-2, to=1-3]
            \arrow["{g_1}"', from=1-2, to=2-2]
            \arrow["0"', from=1-3, to=2-3]
            \arrow[from=2-1, to=2-2]
            \arrow[from=2-2, to=2-3]
        \end{tikzcd}
    \end{displaymath}
    is commutative since the compositions $E \xrightarrow {g_1} E \to C_1$ and $C_1 \to E[1] \xrightarrow {g_1} E[1]$ are zero.
    Continuing, setting $C_{i + 1}$ equal to the cone of the map $g_ i \colon C_ i \to C_ i$, we obtain a perfect complex $C_ s\in \operatorname{Perf}(X)\cap D^{\leq 0}_{\operatorname{qc},T}(X)$ whose restriction to $U$ gives 
    \begin{displaymath}
        F \oplus F[1]^{\oplus s} \oplus F[2]^{\oplus {\binom{s}{2}}} \oplus \ldots \oplus F[s] .
    \end{displaymath}
    By \Cref{lem:lift_roofs}, we can find morphisms of perfect complexes $\beta \colon C^\prime \to C_ s$ and $\gamma \colon C^\prime \to C_ s$ in $D^{\leq 0}_{\operatorname{qc},T}(X)$  such that $\mathbf{L}j^\ast \beta$ is an isomorphism and such that $\mathbf{L}j^\ast \gamma \circ (\mathbf{L}j^\ast \beta)^{-1}$ is the morphism 
    \begin{displaymath}
        F \oplus F[1]^{\oplus s} \oplus F[2]^{\oplus {\binom{s}{2}}} \oplus \ldots \oplus F[s] \to F \oplus F[1]^{\oplus s} \oplus F[2]^{\oplus {\binom{s}{2}}} \oplus \ldots \oplus F[s] 
    \end{displaymath}
    which is the identity on all summands except for $F$ where it is zero. By \Cref{lem:lift_roofs}, we also have 
    \begin{displaymath}
        C^\prime = C_ s \otimes ^\mathbf {L} I\in \operatorname{Perf}(X) \cap D^{\leq 0}_{\operatorname{qc},T}(X) 
    \end{displaymath}
    for some $I\in \operatorname{Perf}(X) \cap D^{\leq 0}_{\operatorname{qc}}(X)$ (see proof of \Cref{lem:lift_roofs}). 
    Hence, the nullity of $g_ j^2 \circ 1_{C_ s}$ implies the same thing for $C^\prime$. 
    Thus, $C^\prime$ is supported on $T$ as well.
    Then $\text{Cone}(\gamma )$ yields the desired object, which belongs to $D^{\leq 0}_{\operatorname{qc},T}(X)$.
\end{proof}

\begin{lemma}
    \label{lem:almost_lift_psuedoapprox}
    Let $X$ be a quasi-compact quasi-separated scheme, $u\colon U \to X$ be a quasi-compact open immersion, $T\subseteq |X|$ be closed with $|X|\setminus T$ retrocompact, and $E\in D^{\leq 0}_{\operatorname{qc}}(X)$.
    For every $\alpha \colon P \to \mathbf{L}u^\ast E$ with $P\in \operatorname{Perf}(U)\cap D^{\leq 0}_{\operatorname{qc},T\cap |U|}(U)$, there exists $R\in \operatorname{Perf}(X)\cap D^{\leq 0}_{\operatorname{qc},T}(X)$ and $\beta \colon R \to E$ such that $P$ is a direct summand of $\mathbf{L}u^\ast R$, compatibly with $\alpha$ and $\mathbf{L}u^\ast \beta$.
\end{lemma}

\begin{proof}
    Since $X$ is quasi-compact, there exists $m$ such that $X = U \cup V_1 \cup \cdots \cup V_ m$ for some affine opens $V_ j$ of $X$. 

    First, say we are in the situation where $X = U \cup V$ with $V$ affine (i.e.\ $m=1$). 
    Form the fibered square 
    \begin{displaymath}
        % https://q.uiver.app/#q=WzAsNCxbMSwwLCJWIl0sWzEsMSwiWCJdLFswLDEsIlUiXSxbMCwwLCJVXFxjYXAgViJdLFswLDEsInYiXSxbMiwxLCJ1IiwyXSxbMywyLCJ2XlxccHJpbWUiLDJdLFszLDAsInVeXFxwcmltZSJdLFszLDEsInQiLDFdXQ==
        \begin{tikzcd}
            {U\cap V} & V \\
            U & X
            \arrow["{u^\prime}", from=1-1, to=1-2]
            \arrow["{v^\prime}"', from=1-1, to=2-1]
            \arrow["t"{description}, from=1-1, to=2-2]
            \arrow["v", from=1-2, to=2-2]
            \arrow["u"', from=2-1, to=2-2]
        \end{tikzcd}
    \end{displaymath}
    where $u,v$ are the associated open immersions. 
    By \Cref{lem:perf_up_to_summand_support}, there exist a perfect object $Q$ in $\operatorname{Perf}(V) \cap D^{\leq 0}_{\operatorname{qc},T \cap V}(V)$ and an isomorphism $\mathbf{L}(u^\prime)^\ast Q \to \mathbf{L}(v^\prime)^\ast (P \oplus P[1])$. 
    Moreover, from \Cref{lem:lift_roofs}, we can replace $Q$ by $Q \otimes ^\mathbf {L} I$ (still supported on $T \cap V$ and in $D^{\leq 0}_{\operatorname{qc}}(V)$, see proof of \Cref{lem:lift_roofs} for details on $I$) and assume that the map 
    \begin{displaymath}
        \mathbf{L}(u^\prime)^\ast Q \to \mathbf{L}(v^\prime)^\ast (P \oplus P[1]) \to \mathbf{L}(v^\prime)^\ast P \to\mathbf{L}t^\ast E 
    \end{displaymath}
    lifts to $Q \to \mathbf{L}v^\ast E$. By \cite[\href{https://stacks.math.columbia.edu/tag/08DG}{Tag 08DG}]{StacksProject}, we can find a morphism $a \colon R \to E$ of $D(X)$ such that $\mathbf{L}u^\ast a$ is isomorphic to $P \oplus P[1] \to \mathbf{L}u^\ast E$ and $\mathbf{L}v^\ast a$ isomorphic to $Q \to \mathbf{L}v^\ast E$. 
    Consequently, $R$ is perfect and supported on $T$. 
    Moreover, $R$ is affine locally in $D^{\leq 0}_{\operatorname{qc}}$, which implies $R\in D^{\leq 0}_{\operatorname{qc}}(X)$. 

    Now for the general case. 
    Set $X_k:= U \cup V_1 \cup \cdots \cup V_k$. 
    As $X$ is quasi-separated, $|U|\setminus (Z \cap |U|) = U\cap (|X|\setminus Z)$ is quasi-compact.
    We can apply the $m=1$ construction above iteratively to $X_{k-1}\subseteq X_k$. 
    At every step, the previously constructed perfect complex remains a direct summand after restriction. 
    Moreover, perfectness, prescribed support, and desired vanishing in cohomology are preserved.
    Thus, we can complete the proof by an inductive argument.
\end{proof}

\begin{remark}
    \label{rmk:psuedo_lifts}
    The following is directly from \cite[\href{https://stacks.math.columbia.edu/tag/09IN}{Tag 09IN}]{StacksProject}.
    Our proof of \Cref{lem:almost_lift_psuedoapprox} shows that
    \begin{displaymath}
        \mathbf{L}u^\ast R \cong  P \oplus P^{\oplus n_1}[1] \oplus \cdots \oplus P^{\oplus n_ m}[m]
    \end{displaymath}
    for some $m \geq 0$ and $n_ j \geq 0$. 
    Hence, the highest degree cohomology sheaf of $\mathbf{L}u^\ast R$ equals that of $P$. 
    By repeating the construction for the map $P^{\oplus n_1}[1] \oplus \ldots \oplus P^{\oplus n_ m}[m] \to \mathbf{L}u^\ast R$, taking cones, and using induction we obtain isomorphisms of cohomology sheaves of $\mathbf{L}u^\ast R$ and $P$ above any given degree.
    More explicitly, \Cref{lem:almost_lift_psuedoapprox} constructs a morphism
    \begin{displaymath}
        a\colon \mathbf{L}u^\ast R \cong P \oplus P^{\oplus n_1}[1] \oplus \cdots \oplus P^{\oplus n_ m}[m],
    \end{displaymath}
    such that the natural morphism 
    \begin{displaymath}
        (\alpha,0,0,\ldots,0) \colon P \oplus P^{\oplus n_1}[1] \oplus \cdots \oplus P^{\oplus n_ m}[m] \to \mathbf{L}u^\ast E
    \end{displaymath}
    factors as $\mathbf{L}u^\ast \beta \circ a^{-1}$.
    Here $(\alpha,0,0,\ldots,0)$ is the morphism induced by $\alpha$ and the zero morphism for the other factors of the direct sum to $\mathbf{L}u^\ast E$. 
\end{remark}

\begin{lemma}
    \label{lem:lisse-etale_pullback_and_qc_pullback}
    Let $f\colon Y\to X$ be a flat morphism of algebraic spaces. 
    If $E\in D_{\operatorname{qc}}(X)$ has finite type cohomology in degree $n$, then $\mathbf{L}f^\ast E$ has finite type cohomology in degree $n$.
\end{lemma}

\begin{proof}
    This follows from flatness of $f$ and \cite[\href{https://stacks.math.columbia.edu/tag/03DO}{Tag 03DO}]{StacksProject}.
    %%NOTE: $\mathcal{H}^n (\mathbf{L}f^\ast E) \cong f^\ast \mathcal{H}^n(E)$ for all $n\in \mathbb{Z}$ via flatness.
\end{proof}

\begin{proposition}
    \label{prop:qcqs_scheme_pseudoapprox}
    Let $X$ be a quasi-compact quasi-separated scheme. 
    Then $X$ satisfies pseudoapproximation by perfect complexes.
\end{proposition}

\begin{proof} 
    We induct on the minimal number $n$ of components for an affine open cover of the scheme. 
    The case $n=1$ is covered by \Cref{lem:quasi-affine_pseudoapprox}. 
    Assume we have proved the claim for all quasi-compact quasi-separated schemes with an affine open cover consisting of at most $n$ components. 
    Let $X=\cup^{n+1}_{i=1} U_i$ be a minimal length affine open cover.
    Denote by $s_i \colon U_i \to X$ the associated open immersion of each $i$.
    Choose a pair $(Z, E)$ where $Z\subseteq |X|$ closed with quasi-compact complement and $E\in D^{\leq 0}_{\operatorname{qc},Z}(X)$ such that $\mathcal{H}^0 (E)$ is of finite type.
    Set $U=\cup^n_{i=1}U_i$. 
    Define $s\colon U \to X$ to be the associated open immersion of $U$.

    By the induction hypothesis, $U$ satisfies pseudoapproximation by perfect complexes. 
    Moreover, from \cite[Lemma A.2]{GuisadoVillaalgordo/Lank:2026} and \Cref{lem:lisse-etale_pullback_and_qc_pullback}, $\mathbf{L}s^\ast E \in D^{\leq 0}_{\operatorname{qc},Z\cap |U|}(U)$ and $\mathcal{H}^0 (\mathbf{L}s^\ast E)$ is of finite type. 
    There exist $P\in \operatorname{Perf}(U) \cap D^{\leq 0}_{\operatorname{qc},Z\cap |U|}(U)$ and $\alpha \colon P\to \mathbf{L}s^\ast E$ which induces a surjective morphism in degree zero. 
    Applying \Cref{lem:almost_lift_psuedoapprox}, we can find $R_1 \in \operatorname{Perf}(X) \cap D^{\leq 0}_{\operatorname{qc},Z}(X)$ and $\alpha_1 \colon R_1 \to E$ such that $P$ is a direct summand of $\mathbf{L}s^\ast R_1$ compatible with $\alpha$ and $\mathbf{L}s^\ast \alpha_1$.
    It is easy to see from this compatibility that $\mathcal{H}^0 (\mathbf{L}s^\ast \alpha_1)$ is surjective. 

    Since $U_{n+1}$ is affine, \Cref{lem:quasi-affine_pseudoapprox} implies it satisfies pseudoapproximation by perfect complexes. 
    Hence, by similar reasoning above, we can find $R_2 \in \operatorname{Perf}(X) \cap D^{\leq 0}_{\operatorname{qc},Z}(X)$ and $\alpha_2 \colon R_2 \to E$ such that $\mathcal{H}^0 (\mathbf{L}s_{n+1}^\ast \alpha_2)$ is surjective. 
    Now, take the natural morphism $\alpha^\prime \colon R_1 \oplus R_2 \to E$ induced by the $\alpha_i \colon R_i \to E$. 
    Then $\mathcal{H}^0 (\alpha^\prime)$ is a surjective morphism, which can be checked after restriction to both components of the open cover $U \cup U_{n+1}=X$. 
    Moreover, $R_1 \oplus R_2$ is in $\operatorname{Perf}(X) \cap D^{\leq 0}_{\operatorname{qc},Z}(X)$, and so we obtain the desired pseudoapproximation. 
\end{proof}

\begin{proof}
    [Proof of \Cref{cor:psuedoapprox}]
    This is immediate from \Cref{prop:approx_by_comorphismcts_implies_Thomason,prop:qcqs_scheme_pseudoapprox}.
\end{proof}

% \begin{remark}
%     There is a generalization which might appeal to some readers. 
%     Instead, the \Cref{def:approximation} could be more restrictive by asking for $E\in D^{\leq 0}_{\operatorname{qc},Z}(X)$ to satisfy $\mathcal{H}^i(E)$ being of finite type with $i\in [b,0]$ for some $b\in [-\infty,0]\cap \mathbb{Z}$. 
%     This could be called \textit{$b$-pseudoapproximation by perfect complexes}.  
%     To check that affine schemes satisfy $b$-pseudoapproximation by perfect complexes, one can induct with \Cref{lem:quasi-affine_pseudoapprox} and taking direct sums of shifts of Koszul complexes. 
%     Then the rest of our results above go through to show all quasi-compact quasi-separated schemes are satisfy $b$-pseudoapproximation by perfect complexes.
%     However, this notion goes beyond the scope of our needs, and we do not study this any further.
% \end{remark}

%%%%%%%%%%%%%%%%%%%%%%%%%%%%%%%%%%%%%
\section{Standard aisle}
\label{sec:standard_aisle}
%%%%%%%%%%%%%%%%%%%%%%%%%%%%%%%%%%%%%

We prove compact generation of the standard aisles. 

\begin{lemma}
    \label{lem:cpt_gen_aisles}
    Let $\mathcal{T}$ be a compactly generated triangulated category. 
    Suppose $\mathcal{U}$ is a cocomplete preaisle on $\mathcal{T}$. 
    If there exists $\mathcal{P}\subseteq \mathcal{T}^c \cap \mathcal{U}$ such that every nonzero $E\in \mathcal{U}$ admits a nonzero morphism $P\to E$ with $P\in \mathcal{P}$, then $\mathcal{U} = \overline{\langle \mathcal{P}\rangle}^{(-\infty,0]}$. 
\end{lemma}

\begin{proof}
    By \cite[Theorem 2.3]{Neeman:2021}, $\overline{\langle \mathcal{P}\rangle}^{(-\infty,0]}$ is an aisle. 
    By construction, we can see that $\overline{\langle \mathcal{P}\rangle}^{(-\infty,0]} \subseteq \mathcal{U}$.
    We check the reverse inclusion. 
    Choose $E\in \mathcal{U}$. 
    Consider the truncation triangle of $E$,
    \begin{displaymath}
        \tau^{\leq 0} E \to E \to \tau^{\geq 1} E \to (\tau^{\leq 0} E)[1]
    \end{displaymath}
    with respect to $\overline{\langle \mathcal{P}\rangle}^{(-\infty,0]}$. 
    Since $\tau^{\leq 0} E$ and $E$ belong to $\mathcal{U}$, it follows that $\tau^{\geq 1} E \in \mathcal{U}$. 
    The hypothesis asserts there exists $P\in \mathcal{P}$ and a nonzero morphism $P \to \tau^{\geq 1} E$. 
    However, this is absurd since $\tau^{\geq 1} E\in \{ Q [n] :n \geq 0, Q\in \mathcal{P}\}^\perp$.
    Thus, $\tau^{\leq 0} E \to E$ is an isomorphism, which gives the claim.
\end{proof}

\begin{lemma}
    \label{lem:support_for_perfect}
    Let $X$ be a quasi-compact quasi-separated algebraic space, $P\in \operatorname{Perf}(X)$, and $s\in |X|$. The following are equivalent:
    \begin{enumerate}
        \item $s\in \operatorname{supp}(P)$
        \item $\mathbf{L}i^\ast P \not\cong 0$ for every representative $i\colon \operatorname{Spec}(k)\to X$ of $p$.
        \item there exists a representative $i\colon \operatorname{Spec}(k)\to X$ of $p$ such that $\mathbf{L}i^\ast P \not\cong 0$.
    \end{enumerate}
\end{lemma}

\begin{proof}
    Choose an \'{e}tale presentation $s\colon U \to X$ be an affine scheme. 
    Note that $\mathbf{L}s^\ast P$ is perfect.
    By flatness of $s$, we know that $\mathcal{H}^n (\mathbf{L}s^\ast P)\cong s^\ast \mathcal{H}^n (P)$.
    Thus, applying \cite[\href{https://stacks.math.columbia.edu/tag/07TY}{Tag 07TY}]{StacksProject}, it suffices to prove the claim for $\mathbf{L}s^\ast P$.
    This follows from \cite[Lemma 3.3]{Thomason:1997}.
\end{proof}

\begin{lemma}
    \label{lem:nbhd_cpt_gen_aisle}
    Consider an elementary distinguished square
    \begin{displaymath}
        % https://q.uiver.app/#q=WzAsNCxbMSwwLCJYXlxccHJpbWUiXSxbMSwxLCJYIl0sWzAsMSwiVSJdLFswLDAsIlVeXFxwcmltZSJdLFswLDEsImYiXSxbMiwxLCJqIiwyXSxbMywyLCJmXlxccHJpbWUiLDJdLFszLDAsImpeXFxwcmltZSJdXQ==
        \begin{tikzcd}
            {U^\prime} & {X^\prime} \\
            U & X
            \arrow["{j^\prime}", from=1-1, to=1-2]
            \arrow["{f^\prime}"', from=1-1, to=2-1]
            \arrow["f", from=1-2, to=2-2]
            \arrow["j"', from=2-1, to=2-2]
        \end{tikzcd}
    \end{displaymath}
    where $j$ is a quasi-compact open immersion and $f$ is an \'{e}tale presentation. 
    Set $Z:= |X|\setminus |U|$. 
    Assume for all closed subsets $S\subseteq |U|$ with quasi-compact complement,
    \begin{displaymath}
        D^{\leq 0}_{\operatorname{qc},S}(U) 
        = \overline{\langle \operatorname{Perf} (U)
        \cap
        D_{\operatorname{qc},S}^{\leq 0}(U) \rangle}^{(-\infty,0]}.
    \end{displaymath}
    Then for all closed subsets $T\subseteq |X|$ with quasi-compact complement,
    \begin{displaymath}
        D^{\leq 0}_{\operatorname{qc},T}(X) 
        = \overline{\langle \operatorname{Perf} (X)
        \cap
        D_{\operatorname{qc},T}^{\leq 0}(X) \rangle}^{(-\infty,0]}.
    \end{displaymath}
\end{lemma}

\begin{proof}
    We can impose $T$ be nonempty.
    Fix a subset $T \subseteq |X|$ which is closed with quasi-compact complement and nonzero $E\in D_{\operatorname{qc},T}^{\leq 0}(X)$.
    Assume that $D^{\leq 0}_{\operatorname{qc},T\cap |U|}(U)$ is compactly generated by 
    \begin{displaymath}
        \operatorname{Perf}(U)
        \cap
        D_{\operatorname{qc},T\cap |U|}^{\leq 0}(U).
    \end{displaymath}
    By \Cref{prop:qcqs_scheme_pseudoapprox,prop:approx_by_comorphismcts_implies_Thomason}, $D^{\leq 0}_{\operatorname{qc},f^{-1}(T\cap Z)}(X^\prime)$ is compactly generated by 
    \begin{displaymath}
        \operatorname{Perf}(X^\prime)
        \cap
        D_{\operatorname{qc},f^{-1}(T\cap Z)}^{\leq 0}(X^\prime).
    \end{displaymath}
    
    Consider the case $\operatorname{supp}(E)\subseteq Z$. 
    By \cite[Lemma A.6]{GuisadoVillaalgordo/Lank:2026}, $\mathbf{L}f^\ast$ and $\mathbf{R}f_\ast$ yield a $t$-exact adjoint equivalence 
    \begin{displaymath}
        \mathbf{L}f^\ast \colon D_{\operatorname{qc},Z}(X) \leftrightarrows D_{\operatorname{qc},f^{-1}(Z)}(X^\prime) \colon \mathbf{R}f_\ast.
    \end{displaymath}
    By \Cref{prop:qcqs_scheme_pseudoapprox}, $X^\prime$ satisfies pseudoapproximation, and hence we can appeal to \Cref{prop:approx_by_comorphismcts_implies_Thomason}.
    Since $\mathbf{L}f^\ast E \in D_{\operatorname{qc},f^{-1}(Z)}(X^\prime)$, there exist 
    \begin{displaymath}
        Q\in 
        \operatorname{Perf}(X^\prime)
        \cap
        D_{\operatorname{qc},f^{-1}(T\cap Z)}^{\leq 0}(X^\prime).
    \end{displaymath}
    and a nonzero morphism $Q \to \mathbf{L}f^\ast E$.
    Then $\mathbf{R}f_\ast Q$ is perfect complex which is supported on $T\cap Z$. 
    Note that the unit $E \to \mathbf{R}f_\ast \mathbf{L}f^\ast E$ is an isomorphism.
    Composing with inverse of the unit yields 
    \begin{displaymath}
        \mathbf{R}f_\ast Q \to \mathbf{R}f_\ast \mathbf{L}f^\ast E \to E.
    \end{displaymath}
    This morphism is nonzero, and $t$-exactness implies $\mathbf{R}f_\ast Q \in D^{\leq 0}_{\operatorname{qc},T}(X)$. 

    Now, we prove the general case, e.g.\ where $\mathbf{L}j^\ast E\cong 0$. 
    By \cite[Lemma A.2]{GuisadoVillaalgordo/Lank:2026} and flatness of $j$, one has $\mathbf{L}j^\ast E \in D_{\operatorname{qc},T\cap |U|}^{\leq 0}(U)$. 
    Applying our hypothesis on $U$, there exist
    \begin{displaymath}
        P_U
        \in
        \operatorname{Perf}(U)
        \cap
        D_{\operatorname{qc},T\cap |U|}^{\leq 0}(U)
    \end{displaymath}
    and a nonzero morphism $u\colon P_U \to \mathbf{L}j^\ast E$. 
    Moreover, we obtain a morphism,
    \begin{displaymath}
        \theta \colon \mathbf{L}(f^\prime)^\ast P_U 
        \xrightarrow{\mathbf{L}(f^\prime)^\ast u} \mathbf{L}(f^\prime)^\ast  \mathbf{L}j^\ast E
        \cong \mathbf{L}(j^\prime)^\ast  \mathbf{L}f^\ast E.
    \end{displaymath}

    Applying \Cref{lem:almost_lift_psuedoapprox} and \Cref{rmk:psuedo_lifts} to $\theta$, there exists $R\in \operatorname{Perf}(X^\prime)\cap D^{\leq 0}_{\operatorname{qc},f^{-1}(T)}(X^\prime)$, a morphism $b \colon R \to \mathbf{L}f^\ast E$, and an isomorphism
    \begin{displaymath}
        \beta \colon \mathbf{L}(j^\prime)^\ast R 
        \cong \mathbf{L}(f^\prime)^\ast P_U \oplus \mathbf{L}(f^\prime)^\ast P_U^{\oplus n_1}[1] \oplus \ldots \oplus \mathbf{L}(f^\prime)^\ast P_U^{\oplus n_ m}[m],
    \end{displaymath}
    for some $m \geq 0$ and $n_ j \geq 0$ such that the natural morphism 
    \begin{displaymath}
        a^\prime \colon \mathbf{L}(f^\prime)^\ast P_U \oplus \mathbf{L}(f^\prime)^\ast P_U^{\oplus n_1}[1] \oplus \cdots \oplus \mathbf{L}(f^\prime)^\ast P_U^{\oplus n_ m}[m] \to \mathbf{L}(j^\prime)^\ast \mathbf{L}f^\ast E
    \end{displaymath}
    factors as $\mathbf{L}j^\ast b \circ \beta^{-1}$.
    Here, $\mathbf{L}(f^\prime)^\ast P_U$ is a direct summand of $\mathbf{L}(j^\prime)^\ast R$, compatibly with $\theta$ and $\mathbf{L}(j^\prime)^\ast b$.
    Set $A$ to be 
    \begin{displaymath}
        P_U \oplus P_U^{\oplus n_1}[1] \oplus \ldots \oplus P_U^{\oplus n_ m}[m].
    \end{displaymath}
    Let
    \begin{displaymath}
        a\colon
        A \to P_U
        \xrightarrow{u}
        \mathbf{L}j^\ast E
    \end{displaymath}
    be the morphism induced by projection onto the unshifted
    $P_U$-summand. 
    Then there exists a commutative diagram
    \begin{displaymath}
        % https://q.uiver.app/#q=WzAsNCxbMCwxLCJcXG1hdGhiZntMfShmXlxccHJpbWUpXlxcYXN0IFBfVSBcXG9wbHVzIFxcbWF0aGJme0x9KGZeXFxwcmltZSleXFxhc3QgUF9VXntcXG9wbHVzIG5fMX1bMV0gXFxvcGx1cyBcXGNkb3RzIFxcb3BsdXMgXFxtYXRoYmZ7TH0oZl5cXHByaW1lKV5cXGFzdCBQX1Vee1xcb3BsdXMgbl8gbX1bbV0iXSxbMCwwLCJcXG1hdGhiZntMfShmXlxccHJpbWUpXlxcYXN0IChQX1UgXFxvcGx1cyBQX1Vee1xcb3BsdXMgbl8xfVsxXSBcXG9wbHVzIFxcY2RvdHMgXFxvcGx1cyBQX1Vee1xcb3BsdXMgbl8gbX1bbV0pIl0sWzIsMSwiXFxtYXRoYmZ7TH0oal5cXHByaW1lKV5cXGFzdCBcXG1hdGhiZntMfWZeXFxhc3QgRS4iXSxbMiwwLCJcXG1hdGhiZntMfShmXlxccHJpbWUpXlxcYXN0IFxcbWF0aGJme0x9al5cXGFzdCBFIl0sWzEsMywiXFxtYXRoYmZ7TH0oZl5cXHByaW1lKV5cXGFzdCBhIl0sWzAsMiwiYV5cXHByaW1lIl0sWzEsMCwiXFxjb25nIiwyXSxbMywyLCJcXGNvbmciLDJdXQ==
        \begin{tikzcd}
            {\mathbf{L}(f^\prime)^\ast (A)} && {\mathbf{L}(f^\prime)^\ast \mathbf{L}j^\ast E} \\
            {\mathbf{L}(f^\prime)^\ast P_U \oplus \mathbf{L}(f^\prime)^\ast P_U^{\oplus n_1}[1] \oplus \cdots \oplus \mathbf{L}(f^\prime)^\ast P_U^{\oplus n_ m}[m]} && {\mathbf{L}(j^\prime)^\ast \mathbf{L}f^\ast E.}
            \arrow["{\mathbf{L}(f^\prime)^\ast a}", from=1-1, to=1-3]
            \arrow["="', from=1-1, to=2-1]
            \arrow["\cong"', from=1-3, to=2-3]
            \arrow["{a^\prime}", from=2-1, to=2-3]
        \end{tikzcd}
    \end{displaymath}

    We want to apply \cite[\href{https://stacks.math.columbia.edu/tag/08HB}
    {Tag 08HB}]{StacksProject}. 
    Define an overlap isomorphism
    \begin{displaymath}
        \begin{aligned}
            c\colon
            \mathbf{L}(f^\prime)^\ast A
            &\xrightarrow{\ \beta^{-1}\ }
            \mathbf{L}(j^\prime)^\ast R.
        \end{aligned}
    \end{displaymath}
    By
    \cite[\href{https://stacks.math.columbia.edu/tag/08HB}
    {Tag 08HB}]{StacksProject},
    there exist an object $P\in D(X)$, and isomorphisms $F\colon \mathbf{L}j^\ast P \to A$ and $G\colon \mathbf{L}f^\ast P \to R$ such that $c = \mathbf{L}(j^\prime)^\ast G \circ (\mathbf{L}(f^\prime)^\ast F)^{-1}$. 

    The local morphisms are compatible on $U^\prime$. Indeed,
    \begin{displaymath}
        \begin{aligned}
            \mathbf{L}(j^\prime)^\ast b\circ c
            &= \mathbf{L}(j^\prime)^\ast b \circ \beta^{-1} 
            \\&= \mathbf{L}(f^\prime)^\ast a \circ \beta\circ\beta^{-1} 
            \\&= \mathbf{L}(f^\prime)^\ast a.
        \end{aligned}
    \end{displaymath}
    Consequently,
    \begin{displaymath}
        \mathbf{L}(j^\prime)^\ast(b\circ G)
        =
        \mathbf{L}(f^\prime)^\ast(a\circ F).
    \end{displaymath}
    Then
    \cite[\href{https://stacks.math.columbia.edu/tag/08HB}
    {Tag 08HB}]{StacksProject} yields a morphism $\phi\colon P\to E$ such that
    $\mathbf{L}j^\ast \phi =a\circ F$ and $\mathbf{L}f^\ast \phi =b\circ G$. 
    Since $f$ is an \'{e}tale surjective morphism, $\mathbf{L}f^\ast \phi =b\circ G$ and the isomorphism $G\colon \mathbf{L}f^\ast P \to R$ implies that $\phi$ is nonzero and that 
    \begin{displaymath}
        P\in \operatorname{Perf}(X) \cap D^{\leq 0}_{\operatorname{qc},T}(X).
    \end{displaymath}
    Indeed, the latter claim follows from \cite[Lemma 3.5]{GuisadoVillaalgordo/Lank:2026} for $P$ being perfect on $X$, \cite[Lemma 2.10]{GuisadoVillaalgordo/Lank:2026} to check that $P\in D^{\leq 0}_{\operatorname{qc}}(X)$, and applying \Cref{lem:support_for_perfect} to verify that $\operatorname{supp}(P)=Z$.
    To see that $\phi$ is not zero, use the fact that $u$ is not zero and $\mathbf{L}j^\ast \phi = a \circ F$. 
    Indeed, $a$ is nonzero because its restriction to the $P_U$-summand is $u\not=0$. 
    Therefore, from \Cref{lem:cpt_gen_aisles}, this completes the proof.
\end{proof}

The following is an equivalent reformulation of \Cref{prop:psuedoapprox}.

\begin{proposition}
    \label{prop:qcqs_cpt_gen_aisle}
    Let $X$ be a quasi-compact quasi-separated algebraic space. 
    Then for all closed subsets $T\subseteq |X|$ with quasi-compact complement,
    \begin{displaymath}
        D^{\leq 0}_{\operatorname{qc},T}(X) 
        = \overline{\langle \operatorname{Perf} (X)
        \cap
        D_{\operatorname{qc},T}^{\leq 0}(X) \rangle}^{(-\infty,0]}.
    \end{displaymath}
\end{proposition}

\begin{proof}
    By \cite[\href{https://stacks.math.columbia.edu/tag/07ST}
    {Tag 07ST}]{StacksProject}, there exist $n\geq 0$ and open subspaces 
    \begin{displaymath}
        \emptyset =: U_0 \xrightarrow{j_0} U_1 \xrightarrow{j_1} \cdots \xrightarrow{j_n} U_{n+1} =: X
    \end{displaymath}
    with the following property: for each $p$ there exists a quasi-compact separated scheme $V_p$ and \'{e}tale presentations $f_p \colon V_p \to U_p$ such that the natural morphisms $f^{-1}_{p+1} (T_{p+1}) \to T_{p+1}$ is an isomorphism where $T_{p+1} := (|U_{p+1}|\setminus |U_p|)_{red}$. 
    In particular, for each $p$, we have an elementary distinguished square $(j_p \colon U_p \to U_{p+1},f_{p+1} \colon V_{p+1} \to U_{p+1})$.
    We will induct on $p$ to show that for all closed subsets $T\subseteq |U_p|$ with quasi-compact complement,
    \begin{displaymath}
        D^{\leq 0}_{\operatorname{qc},T}(U_p) 
        = \overline{\langle \operatorname{Perf}_T (U_p)
        \cap
        D_{\operatorname{qc},T}^{\leq 0}(U_p) \rangle}^{(-\infty,0]}.
    \end{displaymath}
    By \Cref{lem:quasi-affine_pseudoapprox}, the claim holds for $U_0$ because the empty algebraic space is affine. 
    However, the induction step follows by \Cref{lem:nbhd_cpt_gen_aisle}.
    This completes the proof.
\end{proof}

%%%%%%%%%%%%%%%%%%%%%%%%%%%%%%%%%%%%%
\section{Classification}
\label{sec:classification}
%%%%%%%%%%%%%%%%%%%%%%%%%%%%%%%%%%%%%

We classify compactly generated tensor $t$-structures on suitable algebraic spaces.

\begin{definition}
    \label{def:graded_filtration}
    Let $X$ be an algebraic space. 
    We say a Thomason filtration $\phi$ on $X$ is \textbf{supported by perfects} if there exists $\mathcal{S}\subseteq \operatorname{Perf}(X)$ such that for all $n\in \mathbb{Z}$,
    \begin{displaymath}
        \phi(n)=\bigcup_{\substack{j\geq n\\ S\in \mathcal{S}}} \operatorname{supp}(\mathcal{H}^j (S)).
    \end{displaymath}
\end{definition}

\begin{lemma}
    \label{lem:quasi-affine_existence_with_support}
    Let $X=\operatorname{Spec}(R)$ be an affine scheme. 
    For any closed subset $Z$ of $|X|$ with quasi-compact complement, there exists $P\in \operatorname{Perf}(X) \cap D^{\leq 0}_{\operatorname{qc},Z}(X)$ such that $\operatorname{supp}(\mathcal{H}^0 (P))=Z$.
    In particular, $\operatorname{supp}(P)=Z$. 
\end{lemma}

\begin{proof}
    Applying \cite[\href{https://stacks.math.columbia.edu/tag/071Q}{Tag 071Q}]{StacksProject}, we can transfer the problem to the small Zariski site of $X$.
    Fix a closed subset $Z\subseteq |X|$ with quasi-compact complement $U$. 
    Let $j\colon U \to X$ be the open immersion with $|U|:= |X|\setminus Z$. 
    By \cite[\href{https://stacks.math.columbia.edu/tag/0855}{Tag 0855}]{StacksProject}, there exists a closed immersion $i\colon \mathcal{Z} \to X$ which identifies $|\mathcal{Z}|$ with $Z$. 
    It follows that $E := \mathbf{R} i_\ast \mathcal{O}_\mathcal{Z}$ is a quasi-coherent sheaf over $X$ of finite presentation whose support is precisely $Z$. 
    Since $\mathbf{R}i_\ast$ is $t$-exact, $E \in D^{\leq 0}_{\operatorname{qc},Z}(X)$. As pseudoapproximation by perfect complexes holds for $X$ by \Cref{lem:quasi-affine_pseudoapprox}, there is $P \in \operatorname{Perf}(X) \cap D^{\leq 0}_{\operatorname{qc},Z}(X)$ and $P \to E$ such that $\mathcal{H}^0(P) \to E$ is surjective. This yields the following chain of inclusions

    \begin{displaymath}
        Z = \operatorname{supp}(E) \subseteq \operatorname{supp}(\mathcal{H}^0(P)) \subseteq \operatorname{supp}(P) \subseteq Z
    \end{displaymath}
    which establishes the claim.

    % Choose $E\in D^{\leq 0}_{\operatorname{qc},Z}(X)$ such that $\mathcal{H}^0 (E)$ is of finite type. 
    % Let $j\colon U \to X$ be the open immersion with $|U|:= |X|\setminus Z$. 
    % Since $E\in D^{\leq 0}_{\operatorname{qc},Z}(X)$, one has $\mathbf{L}j^\ast \mathbf{R} i_\ast \mathcal{O}_\mathcal{Z} \cong 0$. 

    % Given that $U$ is quasi-compact, there exist $f_1,\ldots,f_r \in R$ such that $U=\cup^r_{i=1} D(f_i)$.
    % Because $X$ is affine and $\mathcal{H}^0 (\mathbf{R} i_\ast \mathcal{O}_\mathcal{Z})$ is of finite type, we can choose global sections $s_1,\ldots,s_c \in H^0 (X,\mathcal{H}^0 (\mathbf{R} i_\ast \mathcal{O}_\mathcal{Z}))$ which generate $H^0 (X,\mathcal{H}^0 (\mathbf{R} i_\ast \mathcal{O}_\mathcal{Z}))$ as an $R$-module. 
    % By \cite[\href{https://stacks.math.columbia.edu/tag/06Z0}{Tag 06Z0}]{StacksProject}, we may regard these as elements of $H^0(X,\mathbf{R} i_\ast \mathcal{O}_\mathcal{Z})$. 
    % For each $a=1,\ldots,c$, \cite[\href{https://stacks.math.columbia.edu/tag/08E3}{Tag 08E3}]{StacksProject} gives an integer $N_a \geq 0$ and a morphism $g_a \colon K_{N_a} \to \mathbf{R} i_\ast \mathcal{O}_\mathcal{Z}$ such that $s_a$ is in the image of the $H^0(X,K_{N_a}) \to H^0(X,\mathbf{R} i_\ast \mathcal{O}_\mathcal{Z})$ where $K_{N_a}$ is the object of $D(X)$ corresponding to the Koszul complex $K^\bullet (f_1^{N_a},\ldots,f_r^{N_a})$ via the equivalence \cite[\href{https://stacks.math.columbia.edu/tag/06Z0}{Tag 06Z0}]{StacksProject}. 
    % The hypothesis of \cite[\href{https://stacks.math.columbia.edu/tag/08E3}{Tag 08E3}]{StacksProject} holds because $\mathbf{L}j^\ast \mathbf{R} i_\ast \mathcal{O}_\mathcal{Z}\cong 0$. 
    % Set $P:= \oplus^c_{a=1} K_{N_a}$. 
    % Consider the natural morphism $g\colon P\to \mathbf{R} i_\ast \mathcal{O}_\mathcal{Z}$ obtained from the $g_a$.
    % %%NOTE: Use def of coproducts and universal properties

    % Each $K_{N_a}$ is a bounded complex of finite locally free sheaves whose cohomology is concentrated in nonpositive degrees, and its cohomology is supported on $Z$. 
    % The image of $H^0(X,P)\to H^0(X,\mathbf{R} i_\ast \mathcal{O}_\mathcal{Z})$ contains all of the $s_a$, and hence is surjective. 
    % Since $X$ is affine, \cite[\href{https://stacks.math.columbia.edu/tag/06Z0}{Tag 06Z0}]{StacksProject} identifies these groups with the global sections of the corresponding degree zero cohomology sheaves. 
    % Hence, the induced morphism $\mathcal{H}^0 (P) \to \mathcal{H}^0 (\mathbf{R} i_\ast \mathcal{O}_\mathcal{Z})$ is surjective. 
    % Hence, we obtain that $Z\subseteq \operatorname{supp}(P)$.
    % Since $P$ is perfect and $\operatorname{supp}(P) \subseteq Z$, we complete the proof.
    % The last claim follows from the proof. 
\end{proof}

\begin{lemma}
    \label{lem:existence_qcqs_prescribed_support}
    Let $X$ be a quasi-compact quasi-separated algebraic space.
    For any closed subset $Z$ of $|X|$ with quasi-compact complement, there exists $P\in \operatorname{Perf}(X) \cap D^{\leq 0}_{\operatorname{qc,Z}}(X)$ such that $\operatorname{supp}(\mathcal{H}^0 (P))=Z$.
    In particular, $\operatorname{supp}(P)=Z$. 
\end{lemma}

\begin{proof}
    We apply \cite[\href{https://stacks.math.columbia.edu/tag/09IT}{Tag 09IT}]{StacksProject}. 
    We may assume $X$ is not empty. 
    By \cite[\href{https://stacks.math.columbia.edu/tag/06NH}{Tag 06NH}]{StacksProject}, $X$ contains a nonempty subspace $X_0$ which is a scheme. 
    Choose a nonempty affine open subscheme $W$ of $X_0$.  
    Then $W$ satisfies the desired claim via \Cref{lem:quasi-affine_existence_with_support}. 

    We check the remaining condition in \cite[\href{https://stacks.math.columbia.edu/tag/09IT}{Tag 09IT}]{StacksProject}.
    Consider an elementary distinguished square given by a quasi-compact open immersion $j\colon U \to X$ and an \'{e}tale morphism $f\colon X^\prime \to X$ from an affine scheme.
    Assume that $U$ satisfies the desired claim. 
    Form the fibered square 
    \begin{displaymath}
        % https://q.uiver.app/#q=WzAsNCxbMSwwLCJYXlxccHJpbWUiXSxbMSwxLCJYIl0sWzAsMSwiVSJdLFswLDAsIlVeXFxwcmltZSJdLFswLDEsImYiXSxbMiwxLCJqIiwyXSxbMywyLCJmXlxccHJpbWUiLDJdLFszLDAsImpeXFxwcmltZSJdXQ==
        \begin{tikzcd}
            {U^\prime} & {X^\prime} \\
            U & {X.}
            \arrow["{j^\prime}", from=1-1, to=1-2]
            \arrow["{f^\prime}"', from=1-1, to=2-1]
            \arrow["f", from=1-2, to=2-2]
            \arrow["j"', from=2-1, to=2-2]
        \end{tikzcd}
    \end{displaymath}
    Fix a closed subset $T\subseteq |X|$ with quasi-compact complement. 
    As $X$ is quasi-separated, we know that $|U|\setminus (T\cap |U|)$ is quasi-compact. 
    By hypothesis, there exists $P\in \operatorname{Perf}(U)\cap D^{\leq 0}_{\operatorname{qc}}(U)$ such that 
    \begin{displaymath}
        \operatorname{supp}(P) = T\cap |U|
    \end{displaymath}
    and $\operatorname{supp}(\mathcal{H}^0 (P))= T\cap |U|$. 
    Applying \Cref{lem:support_for_perfect}, and using that $f^\prime$ is flat, we see that 
    \begin{displaymath}
        \mathbf{L}(f^\prime)^\ast P\in \operatorname{Perf}(U^\prime)\cap D^{\leq 0}_{\operatorname{qc}}(U^\prime)
    \end{displaymath}
    and  
    \begin{displaymath}
        \operatorname{supp}(\mathbf{L}(f^\prime)^\ast P) =\operatorname{supp}(\mathbf{L}(f^\prime)^\ast \mathcal{H}^0 (P)) = f^{-1}(T)\cap |U^\prime|.
    \end{displaymath}
    By \Cref{lem:perf_up_to_summand_support}, there exists 
    \begin{displaymath}
        Q_1 \in \operatorname{Perf}(X^\prime) \cap D^{\leq 0}_{\operatorname{qc},f^{-1}(T)}(X^\prime)
    \end{displaymath}
    such that 
    \begin{displaymath}
        \mathbf{L}(j^\prime)^\ast Q_1 \cong \mathbf{L}(f^\prime)^\ast P \oplus \mathbf{L}(f^\prime)^\ast P[1].
    \end{displaymath}
    Observe, by \Cref{lem:support_for_perfect}, that we only know that
    \begin{displaymath}
        f^{-1}(T) \cap |U^\prime| \subseteq \operatorname{supp}(\mathcal{H}^0 (Q_1)).
    \end{displaymath}
    By \Cref{lem:quasi-affine_existence_with_support}, there exists $Q_2\in D^{\leq 0}_{\operatorname{qc}}(X^\prime)$ such that
    \begin{displaymath}
        \operatorname{supp}(\mathcal{H}^0 (Q_2))= \operatorname{supp}(Q_2)=f^{-1}(T) \cap (|X^\prime|\setminus |U^\prime|).
    \end{displaymath}
    Set $Q= Q_1 \oplus Q_2$. 
    Then $Q\in\operatorname{Perf}(X^\prime)\cap D^{\leq 0}_{\operatorname{qc},f^{-1}(T)}$ and $\operatorname{supp}(\mathcal{H}^0 (Q))=f^{-1}(T)$.
    Moreover, $\mathbf{L}(j^\prime)^\ast Q \cong \mathbf{L}(j^\prime)^\ast Q_1$. 
    Consequently, we obtain an isomorphism 
    \begin{displaymath}
        c\colon \mathbf{L}(f^\prime)^\ast (P \oplus P[1]) 
        \cong \mathbf{L}(f^\prime)^\ast P \oplus \mathbf{L}(f^\prime)^\ast P[1] 
        \cong \mathbf{L}(j^\prime)^\ast Q .
    \end{displaymath}
    By \cite[\href{https://stacks.math.columbia.edu/tag/08HB}{Tag 08HB}]{StacksProject}, there exist an object $R\in D(X)$, and isomorphisms $F\colon \mathbf{L}j^\ast R \to P \oplus P[1]$ and $G\colon \mathbf{L}f^\ast R \to Q$ such that $c = \mathbf{L}(j^\prime)^\ast G \circ (\mathbf{L}(f^\prime)^\ast F)^{-1}$. 

    Denote by $s:= (j,f)\colon U \sqcup X^\prime \to X$ the natural morphism obtained by $j$ and $f$ (see \cite[\href{https://stacks.math.columbia.edu/tag/02WQ}{Tags 02WQ} \& \href{https://stacks.math.columbia.edu/tag/000R}
    {000R}]{StacksProject}). 
    Note that $s$ is \'{e}tale and surjective. 
    Let $p_{\#} \colon \# \to U \sqcup X^\prime$ be the associated open immersions for $\# \in \{U,X^\prime\}$. 
    By construction, it follows that $\mathbf{L}(s \circ p_U)^\ast R \cong \mathbf{L}j^\ast R$ and $\mathbf{L}(s \circ p_{X^\prime})^\ast R \cong \mathbf{L}f^\ast R$. 
    Hence, $\mathbf{L}s^\ast R$ is a perfect complex on $U \sqcup X^\prime$, which belongs to $D^{\leq 0}_{\operatorname{qc}}(U \sqcup X^\prime)$, and is supported on $(T\cap |U|)\sqcup (f^{-1}(T))$. 
    By \cite[Lemma 3.5]{GuisadoVillaalgordo/Lank:2026}, $R$ is perfect on $X$. 
    Moreover, \cite[Lemma 2.10]{GuisadoVillaalgordo/Lank:2026} implies $R\in D^{\leq 0}_{\operatorname{qc}}(X)$. 
    Applying \Cref{lem:support_for_perfect}, it follows that $\operatorname{supp}(R)=T$.
    This completes the proof.
\end{proof}

\begin{lemma}
    \label{lem:affine_supported_by_perfects}
    Let $X$ be a quasi-compact quasi-separated algebraic space.
    Then every Thomason filtration on $X$ is supported by perfects.
\end{lemma}

\begin{proof}
    For each $n \in \mathbb{Z}$, let us write the Thomason set $\phi(n)$ as a union $\cup_{i \in I_n}Z_i^n$ of closed subsets with quasi-compact complements. 
    Using \Cref{lem:existence_qcqs_prescribed_support}, for each $n$ and $i \in I_n$, there exists $P_{Z^i_n}\in \operatorname{Perf}(X)\cap D^{\leq 0}_{\operatorname{qc},Z^i_n}$ such that $\operatorname{supp}(\mathcal{H}^0 (P_{Z^i_n}))= Z^i_n$. 
    Set 
    \begin{displaymath}
        \mathcal{S} := \{P_{Z_i^n}[-n] \mid n \in \mathbb{Z}, i \in I_n\}.
    \end{displaymath}
    It can be verified that $\mathcal{S}$ witnesses that $\phi$ is supported by perfects.
\end{proof}

\begin{lemma}
    \label{lem:cohomology_vanish_from_aisle}
    Let $X$ be a quasi-compact quasi-separated algebraic space. 
    Suppose $E\in D_{\operatorname{qc}}(X)$ and $\mathcal{S}\subseteq  D_{\operatorname{qc}}(X)$. 
    If $E$ is an $n$-fold extension of objects from $\mathcal{S}$, then for any $j\in \mathbb{Z}$ one has 
    \begin{displaymath}
       \operatorname{supp}(\mathcal{H}^j (E))\subseteq \bigcup_{S\in  \mathcal{S}} \operatorname{supp}(\mathcal{H}^j (S)).
    \end{displaymath}
\end{lemma}

\begin{proof}
    Let $U \to X$ be an \'{e}tale presentation. 
    Since $U\to X$ is surjective and flat, the desired support containment may be checked after derived pullback to $U$. 
    Thus, we reduce to the case $X$ is a scheme. 

    We induct on $n$. 
    The base case with $n=1$ is straightforward. 
    Indeed, $E$ is isomorphic to an object in $\mathcal{S}$. 
    Now suppose we have proven the claim for any object that is a $k$-fold extension of objects from $\mathcal{S}$ where $1\leq k \leq  n-1$. 
    
    Let $E$ be an $n$-fold extension of objects from $\mathcal{S}$. 
    This means there is a distinguished triangle
    \begin{displaymath}
        E^\prime \to E \to E^{\prime \prime} \to E^{\prime}[1]
    \end{displaymath}
    where $E^\prime$ is an $(n-1)$-fold extension of objects from $\mathcal{S}$ and $E^{\prime \prime}$ is isomorphic to an object from $\mathcal{S}$. 
    Consider the long exact sequence in cohomology 
    \begin{displaymath}
        \cdots \to \mathcal{H}^j (E^\prime) \to \mathcal{H}^j (E) \to \mathcal{H}^j (E^{\prime \prime}) \to \cdots.
    \end{displaymath}
    Choose 
    \begin{displaymath}
        p \in X\setminus (\cup_{S\in  \mathcal{S}} \operatorname{supp}(\mathcal{H}^j (S))).
    \end{displaymath} 
    Then there exists a long exact sequence in cohomology
    \begin{displaymath}
        \cdots \to \mathcal{H}^j (E^\prime)_p \to \mathcal{H}^j (E)_p \to \mathcal{H}^j (E^{\prime \prime})_p \to \cdots.
    \end{displaymath}
    By the choice of $p$, the induction hypothesis implies 
    \begin{displaymath}
        (\mathcal{H}^j (E^\prime))_p \cong \mathcal{H}^j (E^\prime_p) \cong 0. 
    \end{displaymath} 
    A similar statement holds for $E^{\prime \prime}$. 
    Thus, $p\not \in \operatorname{supp}(\mathcal{H}^j (E))$, completing the proof.
\end{proof}

\begin{lemma}
    \label{lem:thomason_description}
    Let $X=\operatorname{Spec}(R)$ for a commutative ring (which is not necessarily Noetherian). 
    Suppose $\mathcal{P}\subseteq \operatorname{Perf}(X)$. 
    Under the bijection of \cite[Theorem 5.1]{Hrbek:2020}, the corresponding Thomason filtration of $\overline{\langle \mathcal{P}\rangle}^{(-\infty,0]}$ can be also identified with the Thomason filtration on $X$ given by the rule
    \begin{displaymath}
        n \mapsto \bigcup_{\substack{j\geq n \\ P \in \mathcal{P}}} \operatorname{supp}(\mathcal{H}^j (P)).
    \end{displaymath}
\end{lemma}

\begin{proof}
    Applying \cite[\href{https://stacks.math.columbia.edu/tag/071Q}{Tags 071Q}, \href{https://stacks.math.columbia.edu/tag/08H8}{08H8}, \& \href{https://stacks.math.columbia.edu/tag/07TY}{07TY}]{StacksProject}, we can transfer \cite[Theorem 5.1]{Hrbek:2020} to the small Zariski site of $X$ to the small \'{e}tale site. 
    Set $\mathcal{U}:= \overline{\langle \mathcal{P}\rangle}^{(-\infty,0]}$. 
    By \cite[Theorem 5.1]{Hrbek:2020}, the corresponding Thomason filtration of $\mathcal{U}$ is given by the rule
    \begin{displaymath}
        \phi_{\mathcal{U}}\colon k\mapsto \bigcup_{Z_{\mathcal{I}} \textrm{ s.t. } (i_{\mathcal{I}})_\ast \mathcal{O}_{Z_{\mathcal{I}}} [ -k]\in \mathcal{U}} Z_{\mathcal{I}} 
    \end{displaymath}
    where $i_{\mathcal{I}} \colon Z_{\mathcal{I}} \to X$ is a closed subscheme associated to a finitely generated ideal sheaf $\mathcal{I}\subseteq \mathcal{O}_X$. 
    Denote by $\mathcal{D}$ the collection of pairs $(\mathcal{I},-k)$ with $\mathcal{I}$ an ideal sheaf and $k$ any integer appearing above. Fix $n\in \mathbb{Z}$. 
    It suffices to check that 
    \begin{displaymath}
        \bigcup_{\substack{j\geq n \\ P \in \mathcal{P}}} \operatorname{supp}(\mathcal{H}^j (P)) = \bigcup_{Z_{\mathcal{I}} \textrm{ s.t. } (i_{\mathcal{I}})_\ast \mathcal{O}_{Z_{\mathcal{I}}} [ -n]\in \mathcal{U}} Z_{\mathcal{I}}.
    \end{displaymath}

    Let $Z_{\mathcal{I}}$ be a closed subset such that $(i_{\mathcal{I}})_\ast \mathcal{O}_{Z_{\mathcal{I}}} [ -n]\in U$ for some $(\mathcal{I},-n)\in \mathcal{D}$. 
    By \cite[Lemma 5.3(iii)]{Hrbek:2020}, 
    \begin{displaymath}
        \overline{\langle \mathcal{O}_{Z_{\mathcal{I}}} [ -n] \rangle}^{(-\infty,0]} = \overline{\langle K(\mathcal{I}) [ -n] \rangle}^{(-\infty,0]}
    \end{displaymath}
    where $K(\mathcal{I})$ is the Koszul complex on a set of generators for $\mathcal{I}$. 
    By \cite[Lemma 5.1]{Hrbek:2020}, $ K(\mathcal{I})[ -n]$ is a direct summand of an object $A_{\mathcal{I},n}$ which is an $r$-fold extension of objects from the collection $\{ P [t] : t\geq 0, P \in \mathcal{P} \}$.
    Applying \Cref{lem:cohomology_vanish_from_aisle}, it follows that
    \begin{displaymath}
        \begin{aligned}
            Z_{\mathcal{I}} 
            &= \operatorname{supp} (\mathcal{H}^n ( K(\mathcal{I}) [ -n])) 
            \\&\subseteq \operatorname{supp}(\mathcal{H}^n (A_{\mathcal{I},n})) \\&\subseteq \bigcup_{\substack{t\geq 0 \\ P \in \mathcal{P}}} \operatorname{supp}(\mathcal{H}^n (P[t])) 
            \\&\subseteq \bigcup_{\substack{j\geq n \\ P \in \mathcal{P}}} \operatorname{supp}(\mathcal{H}^j (P)).
        \end{aligned}
    \end{displaymath}
    Consequently,
    \begin{displaymath}
        \bigcup_{Z_{\mathcal{I}} \textrm{ s.t. } (i_{\mathcal{I}})_\ast \mathcal{O}_{Z_{\mathcal{I}}} [ -n]\in \mathcal{U}} Z_{\mathcal{I}} \subseteq \bigcup_{\substack{j\geq n \\ P \in \mathcal{P}}} \operatorname{supp}(\mathcal{H}^j (P)).
    \end{displaymath}

    We check the reverse inclusion. 
    Consider the Thomason filtration $\phi_{\mathcal{P}}$ on $X$ given by the rule
    \begin{displaymath}
        n \mapsto \bigcup_{\substack{j\geq n \\ P \in \mathcal{P}}} \operatorname{supp}(\mathcal{H}^j (P)).
    \end{displaymath}
    We want to establish for each $n\in \mathbb{Z}$ that $\phi_{\mathcal{P}}(n) \subseteq \phi_{\mathcal{U}}(n)$. 
    Observe, for each $E\in \mathcal{U}$, that 
    \begin{displaymath}
        \operatorname{supp}(\mathcal{H}^j(E)) \subseteq \phi_{\mathcal{U}}(j)
    \end{displaymath}
    for all $j\in \mathbb{Z}$. 
    Thus, for $P\in \mathcal{P}$ and $j\geq n$, 
    \begin{displaymath}
        \operatorname{supp}(\mathcal{H}^j(P)) \subseteq \phi_{\mathcal{U}}(j) \subseteq \phi_{\mathcal{U}}(n). 
    \end{displaymath}
    Taking the union over all $P\in \mathcal{P}$ gives $\phi_{\mathcal{P}}(n) \subseteq \phi_{\mathcal{U}}(n)$ for all $n\in \mathbb{Z}$.
    %%OLD PROOF:
    % Choose $P\in \mathcal{P}$ and $j\geq n$ such that $\operatorname{supp}(\mathcal{H}^j (P))$ is nonempty. 
    % By \cite[Theorem 5.1]{Hrbek:2020}, 
    % \begin{displaymath}
    %     \overline{ \langle \{ (i_{\mathcal{I}})_\ast \mathcal{O}_{\mathcal{Z}_{\mathcal{I}}} [-k] : (\mathcal{I},-k)\in \mathcal{D}\} \rangle}^{(-\infty,0]} = \mathcal{U}.
    % \end{displaymath}
    % Once more, by \cite[Lemma 5.1]{Hrbek:2020}, $P$ is a direct summand of an object $B$ which is an $r$-fold extension of objects from the collection
    % \begin{displaymath}
    %     \{ (i_{\mathcal{I}})_\ast \mathcal{O}_{\mathcal{Z}_{\mathcal{I}}} [-k+s] : (\mathcal{I},-k)\in \mathcal{D}, s\geq 0\}.
    % \end{displaymath} 
    % Applying \Cref{lem:cohomology_vanish_from_aisle}, we obtain that 
    % \begin{equation}
    %     \label{eq:thomason_description}
    %     \begin{aligned}
    %         \emptyset\not=\operatorname{supp}(\mathcal{H}^j (P))
    %         &\subseteq \operatorname{supp}(\mathcal{H}^j (B))
    %         \\&\subseteq \bigcup_{\substack{(\mathcal{I},-k)\in \mathcal{D} \\ s\geq 0}} \operatorname{supp}(\mathcal{H}^j ((i_{\mathcal{I}})_\ast \mathcal{O}_{\mathcal{Z}_{\mathcal{I}}} [-k+s])).
    %     \end{aligned}
    % \end{equation}
    % Choose any $(\mathcal{I},-k)\in \mathcal{D}$ and $s\geq 0$ above such that 
    % \begin{displaymath}
    %     \operatorname{supp}(\mathcal{H}^j (P)) \cap \operatorname{supp}(\mathcal{H}^j ((i_{\mathcal{I}})_\ast \mathcal{O}_{\mathcal{Z}_{\mathcal{I}}} [-k+s]))\not=\emptyset.
    % \end{displaymath}
    % Note there is at least one such pair from the hypothesis. 
    % This yields
    % \begin{displaymath}
    %     \mathcal{H}^j ((i_{\mathcal{I}})_\ast \mathcal{O}_{\mathcal{Z}_{\mathcal{I}}} [-k+s])\not=0,
    % \end{displaymath}
    % which can only be the case when $j = k -s$. 
    % However, $j\geq n$, and so, $-k + s \leq -n$.
    % Recall that aisles are closed under positive shifts.
    % Since $(i_{\mathcal{I}})_\ast \mathcal{O}_{\mathcal{Z}_{\mathcal{I}}} [-k] \in U$, it follows that $(i_{\mathcal{I}})_\ast \mathcal{O}_{\mathcal{Z}_{\mathcal{I}}} [-n]\in U$. 
    % It is clear that the support of $\mathcal{H}^j ((i_{\mathcal{I}})_\ast \mathcal{O}_{\mathcal{Z}_{\mathcal{I}}} [-k+s])$ is $Z_{\mathcal{I}}$. 
    % Hence, by \eqref{eq:thomason_description}, $\operatorname{supp}(\mathcal{H}^j (P))$ is contained in the union of all such $Z_{\mathcal{I}}$ above. 
    % Therefore,
    % \begin{displaymath}
    %     \operatorname{supp}(\mathcal{H}^j (P)) \subseteq \bigcup_{Z_{\mathcal{I}} \textrm{ s.t. } (i_{\mathcal{I}})_\ast \mathcal{O}_{Z_{\mathcal{I}}} [ -n]\in \mathcal{U}} Z_{\mathcal{I}},
    % \end{displaymath}
    % which completes the proof.
\end{proof}

\begin{remark}
    \label{rmk:proxy_aisle}
    Let $X$ be a quasi-compact quasi-separated algebraic space and $\mathcal{C}\subseteq D_{\operatorname{qc}}(X)$ be a set. 
    By \Cref{prop:psuedoapprox}, $D^{\leq 0}_{\operatorname{qc}}(X)$ is compactly generated by $D^{\leq 0}_{\operatorname{qc}}(X) \cap \operatorname{Perf}(X)$. 
    Then, from \cite[Lemmas 2.6 \& 2.8]{Hrbek/Lank/LeGros/Pavon:2026}, it follows
    \begin{displaymath}
        \overline{\langle \mathcal{C} \rangle}^{(-\infty,0]}_{\otimes} = \overline{\langle (D^{\leq 0}_{\operatorname{qc}}(X) \cap \operatorname{Perf}(X)) \otimes^{\mathbf{L}} \mathcal{C} \rangle}^{(-\infty,0]}.
    \end{displaymath}
\end{remark}

\begin{lemma}
    \label{lem:pullback_t_exactness}
    Let $X$ be a quasi-compact quasi-separated algebraic space. 
    Suppose $f\colon U \to X$ is a smooth morphism from an affine scheme. 
    Then, for any $\mathcal{P}\subseteq \operatorname{Perf}(X)$ such that $\operatorname{preaisle}_{\otimes}(\mathcal{P})$ is an aisle on $D_{\operatorname{qc}}(X)$, one has that $\mathbf{L} f^\ast$ induces a $t$-exact functor
    \begin{displaymath}
        (D_{\operatorname{qc}}(X),\overline{\langle \mathcal{P} \rangle}^{(-\infty,0]}_{\otimes}) \to (D_{\operatorname{qc}}(U),\overline{ \langle \mathbf{L} f^\ast \mathcal{P} \rangle }^{(-\infty,0]}_{\otimes}).
    \end{displaymath}
\end{lemma}

\begin{proof}
    Assume the contrary. 
    Note that all aisles are tensor compatible because $U$ is affine. 
    Hence, by \Cref{rmk:proxy_aisle}, $\mathbf{L}f^\ast \overline{\langle \mathcal{P} \rangle}^{(-\infty,0]}_{\otimes} \subseteq \overline{ \langle \mathbf{L} f^\ast \mathcal{P} \rangle}^{(-\infty,0]}$.
    This shows that $\mathbf{L}f^\ast$ is right $t$-exact. 
    Indeed, from \Cref{lem:dubey_sahoo_corrected}, it suffices to check that 
    \begin{displaymath}
        \mathbf{L}f^\ast ((D^{\leq 0}_{\operatorname{qc}}(X) \cap \operatorname{Perf}(X)) \otimes^{\mathbf{L}} \mathcal{P})
        \subseteq \overline{ \langle \mathbf{L} f^\ast \mathcal{P} \rangle}^{(-\infty,0]} .
    \end{displaymath}
    Since $\mathbf{L}f^\ast (D^{\leq 0}_{\operatorname{qc}}(X) \cap \operatorname{Perf}(X)) \subseteq \overline{\langle \mathcal{O}_U \rangle}^{(-\infty,0]}$, the claim follows. 
    Hence, the assumption ensures $\mathbf{L}f^\ast$ must fail to be left $t$-exact. 
    So, we can find an $E\in (\overline{ \langle\mathcal{P} \rangle}^{(-\infty,0]})^\perp[1]$ such that $\mathbf{L}f^\ast E \not\in (\overline{ \langle \mathbf{L} f^\ast \mathcal{P} \rangle}^{(-\infty,0]})^\perp[1]$. 
    This tells us $\mathbf{L} f^\ast E [-1]\not\in (\overline{\langle \mathbf{L} f^\ast \mathcal{P}\rangle}^{(-\infty,0]}_{\otimes})^\perp$. 
    By \cite[\href{https://stacks.math.columbia.edu/tag/08H6}{Tag 08H6}]{StacksProject}, we know that $\mathbf{L}f^\ast \mathcal{P}\subseteq \operatorname{Perf}(U)$. 
    Moreover, using \cite[Lemma 3.1]{AlonsoTarrio/Lopez/Salorio:2003}, we can find an $A\in \mathcal{P}$ and $j\geq 0$ such that $\operatorname{Hom}(\mathbf{L}f^\ast A[j] , \mathbf{L}f^\ast E[-1])\not=0$. 
    Now we use adjunction to see that,
    \begin{displaymath}
        \begin{aligned}
            \operatorname{Hom} & (\mathbf{L}f^\ast A[j] , \mathbf{L}f^\ast E[-1]) 
            \\&\cong \operatorname{Hom}(\mathcal{O}_U, \operatorname{\mathbf{R}\mathcal{H}\! \mathit{om}}(\mathbf{L}f^\ast A[j], \mathbf{L}f^\ast E[-1]) ) && \textrm{\cite[\href{https://stacks.math.columbia.edu/tag/08DH}{Tag 08DH}]{StacksProject}}
            \\&\cong \operatorname{Hom}(\mathcal{O}_U, \mathbf{L}f^\ast \operatorname{\mathbf{R}\mathcal{H}\! \mathit{om}}(A[j], E[-1]) )
        \end{aligned}
    \end{displaymath}
    where the last isomorphism comes from \cite[Lemmas 2.7 \& 2.11]{GuisadoVillaalgordo/Lank:2026} and \cite[\href{https://stacks.math.columbia.edu/tag/0A8A}{Tag 0A8A}]{StacksProject}.
    This implies that
    \begin{displaymath}
        \mathbf{L}f^\ast \operatorname{\mathbf{R}\mathcal{H}\! \mathit{om}}(A[j], E[-1]) \not\in D^{\geq 1}_{\operatorname{qc}}(U),
    \end{displaymath}
    and because $\mathbf{L}f^\ast$ is $t$-exact with respect to the standard $t$-structures, 
    \begin{displaymath}
        \operatorname{\mathbf{R}\mathcal{H}\! \mathit{om}}(A[j], E[-1]) \not\in D^{\geq 1}_{\operatorname{qc}}(X) = (D^{\leq 0}_{\operatorname{qc}}(X))^\perp.
    \end{displaymath}
    Then we can find a $B\in D^{\leq 0}_{\operatorname{qc}}(X)$ such that
    \begin{displaymath}
        \begin{aligned}
            0 &\not= 
            \operatorname{Hom}(B , \operatorname{\mathbf{R}\mathcal{H}\! \mathit{om}}(A[j], E[-1]) )
            \\&\cong \operatorname{Hom}(B \otimes^{\mathbf{L}} A[j], E[-1] ) && (\textrm{\cite[\href{https://stacks.math.columbia.edu/tag/08DH}{Tag 08DH}]{StacksProject}}).
        \end{aligned}
    \end{displaymath}
    However, $E[-1]$ is in $(\overline{\langle \mathcal{P}\rangle}^{(-\infty,0]}_{\otimes})^\perp$ and $B\otimes^{\mathbf{L}} A[j]\in \overline{\langle \mathcal{P}\rangle}^{(-\infty,0]}_{\otimes}$, which is a contradiction.
\end{proof}

\begin{lemma}
    \label{lem:cumulative-support}
    Let $X$ be a quasi-compact quasi-separated algebraic space,
    $P\in\operatorname{Perf}(X)$, and $n\in\mathbb Z$. Set
    \begin{displaymath}
        \operatorname{supp}^{\geq n}(P)
        :=
        \bigcup_{j\geq n}
        \operatorname{supp}(\mathcal{H}^j(P)).
    \end{displaymath}
    Then $\operatorname{supp}^{\geq n}(P)$ is a Thomason subset of $|X|$.
    Moreover, if $g\colon Y\to X$ is flat, then
    \begin{displaymath}
        g^{-1}(\operatorname{supp}^{\geq n}(P))
        =
        \operatorname{supp}^{\geq n}(\mathbf{L}g^\ast P).
    \end{displaymath}
\end{lemma}

\begin{proof}
    Choose an \'{e}tale presentation $U \to X$. 
    Since $U\to X$ is surjective and flat, the desired properties may be checked after derived pullback to $U$. 
    Thus, we reduce to the case $X$ is a scheme. 

    By the proof of \cite[Lemma 3.3(c)]{Thomason:1997}, it follows that $\operatorname{supp}(\mathcal{H}^j(P))$ is closed. 
    Since $P$ has bounded cohomology, it follows that $\operatorname{supp}^{\geq n}(P)$ is a finite union of closed subsets, and so it is closed.
   Thus, $X$ being quasi-separated implies 
    \begin{displaymath}
        |X| \setminus \bigcup_{j\geq n}
        \operatorname{supp}(\mathcal{H}^j(P)) = \bigcap_{j\geq n} (|X| \setminus \operatorname{supp}(\mathcal{H}^j(P)))
    \end{displaymath}
    is quasi-compact and open in $|X|$. 
    See e.g.\ \cite[\href{https://stacks.math.columbia.edu/tag/054D}{Tag 054D}]{StacksProject}.

    We prove the second claim. 
    This is \'{e}tale local.
    Hence, we can replace the source and target by schemes. 
    Then the result follows from \cite[Lemma 3.3(b)]{Thomason:1997}. 
\end{proof}

\begin{proof}
    [Proof of \Cref{thm:stacky_DB}]
    For a collection $\mathcal P\subseteq\operatorname{Perf}(X)$, set $\mathcal U_{\mathcal P}:=\langle\mathcal P\rangle_{\otimes}^{(-\infty,0]}$ and define
    \[
        \Phi_{\mathcal P}(n)
        :=
        \bigcup_{\substack{P\in\mathcal P\\j\geq n}}
        \operatorname{supp}\bigl(\mathcal{H}^j(P)\bigr).
    \]
    By \Cref{lem:cumulative-support}, it follows that $\Phi_{\mathcal P}$ is a
    Thomason filtration on $X$.

    Choose an \'{e}tale presentation $f\colon U\to X$ with $U$ affine. 
    For a compactly generated tensor aisle
    $\mathcal U_{\mathcal P}$, set
    \begin{displaymath}
        f^\sharp(\mathcal U_{\mathcal P})
        :=
        \mathcal U_{\mathbf{L}f^\ast\mathcal P}
        =
        \left\langle
            \mathbf{L}f^\ast\mathcal P
        \right\rangle_{\otimes}^{(-\infty,0]}.
    \end{displaymath}

    We first check that this construction is independent of the chosen
    generating collection. Suppose $\mathcal U_{\mathcal P} = \mathcal U_{\mathcal Q}$. 
    For every $P\in\mathcal P$, one has $P\in\mathcal U_{\mathcal Q}$. 
    By \Cref{lem:pullback_t_exactness}, the functor
    $\mathbf{L}f^\ast$ is right $t$-exact from $\mathcal U_{\mathcal Q}$ to
    $\mathcal U_{\mathbf{L}f^\ast\mathcal Q}$. 
    Hence, $\mathbf{L}f^\ast P \in \mathcal U_{\mathbf{L}f^\ast\mathcal Q}$ for every $P\in\mathcal P$, and so $\mathcal U_{\mathbf{L}f^\ast\mathcal P} \subseteq \mathcal U_{\mathbf{L}f^\ast\mathcal Q}$. 
    The reverse inclusion follows symmetrically. 
    This shows that $f^\sharp$ is well-defined.

    We next show that $f^\sharp$ is injective. Suppose $\mathcal U_{\mathbf{L}f^\ast\mathcal P} = \mathcal U_{\mathbf{L}f^\ast\mathcal Q}$ where $\mathcal U_{\mathcal P}\not\cong\mathcal U_{\mathcal Q}$. 
    After interchanging $\mathcal P$ and $\mathcal Q$, if necessary,
    choose $E\in \mathcal U_{\mathcal P}\setminus\mathcal U_{\mathcal Q}$. 
    Consider the truncation triangle
    \begin{displaymath}
        \tau_{\mathcal Q}^{\leq0}E
        \to
        E
        \to
        \tau_{\mathcal Q}^{\geq1}E
        \to
        (\tau_{\mathcal Q}^{\leq0}E)[1]
    \end{displaymath}
    with respect to $\mathcal U_{\mathcal Q}$. Since $E\notin\mathcal U_{\mathcal Q}$, one has $\tau_{\mathcal Q}^{\geq1}E\not\cong0$. 
    By \Cref{lem:pullback_t_exactness}, applying $\mathbf{L}f^\ast$ gives the truncation triangle of $\mathbf{L}f^\ast E$ with respect to $\mathcal U_{\mathbf{L}f^\ast\mathcal Q}$. 
    However, 
    \begin{displaymath}
        \mathbf{L}f^\ast E
        \in
        \mathcal U_{\mathbf{L}f^\ast\mathcal P}
        =
        \mathcal U_{\mathbf{L}f^\ast\mathcal Q},
    \end{displaymath}
    and so $\mathbf{L}f^\ast (\tau_{\mathcal Q}^{\geq1}E) \cong 0$. 
    Since $f$ is faithfully flat, $\mathbf{L}f^\ast$ is conservative.
    Thus, we obtain that $\tau_{\mathcal Q}^{\geq1}E=0$, which is a contradiction. 
    Hence, $f^\sharp$ is injective.

    We now define
    \begin{displaymath}
        \Theta\colon
        \left\{
            \begin{array}{c}
            \text{compactly generated tensor aisles}\\
            \text{on }D_{\operatorname{qc}}(X)
            \end{array}
        \right\}
        \to
        \left\{
            \text{Thomason filtrations on }X
        \right\}
    \end{displaymath}
    by $\Theta(\mathcal U_{\mathcal P}):= \Phi_{\mathcal P}$. 

    We check that $\Theta$ is well-defined. Suppose $\mathcal U_{\mathcal P}=\mathcal U_{\mathcal Q}$. 
    By the well-definedness of $f^\sharp$, $\mathcal U_{\mathbf{L}f^\ast\mathcal P} = \mathcal U_{\mathbf{L}f^\ast\mathcal Q}$. 
    Since $U$ is affine, the classification of
    \cite[Theorem 5.1]{Hrbek:2020}, together with \Cref{lem:thomason_description}, gives $\Phi_{\mathbf{L}f^\ast\mathcal P}=\Phi_{\mathbf{L}f^\ast\mathcal Q}$. 
    By \Cref{lem:cumulative-support}, we infer that 
    \begin{displaymath}
         \Phi_{\mathbf{L}f^\ast\mathcal P}
        =
        f^{-1}(\Phi_{\mathcal P})
        \qquad\text{and}\qquad
        \Phi_{\mathbf{L}f^\ast\mathcal Q}
        =
        f^{-1}(\Phi_{\mathcal Q}).
    \end{displaymath}
    Since $f$ is surjective, $\Phi_{\mathcal P} = \Phi_{\mathcal Q}$, and so $\Theta$ is well-defined.

    Next, suppose $\Theta(\mathcal U_{\mathcal P})= \Theta(\mathcal U_{\mathcal Q})$. 
    Then
    \begin{displaymath}
        \Phi_{\mathbf{L}f^\ast\mathcal P}
        =
        f^{-1}(\Phi_{\mathcal P})
        =
        f^{-1}(\Phi_{\mathcal Q})
        =
        \Phi_{\mathbf{L}f^\ast\mathcal Q}.
    \end{displaymath}
    Then \cite[Theorem 5.1]{Hrbek:2020} and \Cref{lem:thomason_description} imply $\mathcal U_{\mathbf{L}f^\ast\mathcal P} = \mathcal U_{\mathbf{L}f^\ast\mathcal Q}$. 
    By the injectivity of $f^\sharp$, it follows that $\mathcal U_{\mathcal P} = \mathcal U_{\mathcal Q}$. 
    Hence, $\Theta$ is injective.

    Finally, let $\phi$ be a Thomason filtration on $X$. 
    By \Cref{lem:affine_supported_by_perfects},
    there exists a collection $\mathcal S_\phi\subseteq\operatorname{Perf}(X)$ such that
    \begin{displaymath}
        \phi(n)
        =
        \bigcup_{\substack{S\in\mathcal S_\phi\\j\geq n}}
        \operatorname{supp}(\mathcal{H}^j(S))
        =
        \Phi_{\mathcal S_\phi}(n)
    \end{displaymath}
    for every $n\in\mathbb Z$. 
    Therefore, 
    \begin{displaymath}
        \Theta
        \left(
            \langle\mathcal S_\phi\rangle_{\otimes}^{(-\infty,0]}
        \right)
        =
        \phi.
    \end{displaymath}
    Consequently, $\Theta$ is surjective and hence bijective.
\end{proof}

%%%%OLD PROOF:

% \begin{lemma}
%     \label{lem:supp_perf_pullback}
%     Let $g\colon Y\to X$ be a morphism quasi-compact quasi-separated algebraic spaces.
%     If $P\in \operatorname{Perf}(X)$, then $g^{-1}\operatorname{supp}(P)=\operatorname{supp}(\mathbf{L}g^\ast P)$.
% \end{lemma}

% \begin{proof}
%     This is \'{e}tale local.
%     Hence, we can replace the source and target by schemes. 
%     Then the result follows from \cite[Lemma 3.3(b)]{Thomason:1997}.
% \end{proof}

% \begin{proof}
%     [Proof of \Cref{thm:stacky_DB}]
%     Let $f\colon U \to X$ be a smooth surjective morphism from an affine scheme. 
%     Set 
%     \begin{displaymath}
%         \mathcal{C}_1 := \{ \textrm{compactly generated } \otimes\!\textrm{-aisle on } D_{\operatorname{qc}}(X)\} 
%     \end{displaymath}
%     and 
%     \begin{displaymath}
%         \mathcal{C}_2 := \{ \textrm{Thomason filtrations on } X\}.
%     \end{displaymath}
%     Define $\Phi\colon \mathcal{C}_1 \to \mathcal{C}_2$ to be the assignment for any $\mathcal{U}\in \mathcal{C}_1$ such that $\mathcal{U}=\overline{\langle \mathcal{P} \rangle}^{(-\infty,0]}_{\otimes}$ for some $\mathcal{P} \subseteq \operatorname{Perf}(X)$, $\Phi(\mathcal{U})$ is the Thomason filtration given by the rule
%     \begin{displaymath}
%         n\mapsto \bigcup_{\substack{j\geq n\\ P\in \mathcal{P}}} \operatorname{supp}(\mathcal{H}^j (P)).
%     \end{displaymath}




%     First, we show there is an injective mapping
%     \begin{equation}
%         \label{eq:classification_map1}
%         \begin{aligned}
%             \{ \otimes\textrm{-aisle on } & D_{\operatorname{qc}}(X) \textrm{ generated by } \mathcal{P} \subseteq \operatorname{Perf}(X) \} 
%             \\&\to \{ \otimes\textrm{-aisle on } D_{\operatorname{qc}}(U) \textrm{ generated by } \mathcal{Q} \subseteq \operatorname{Perf}(U) \}.
%         \end{aligned}
%     \end{equation}
%     This part of the proof is done by contradiction. 
%     Suppose $\mathcal{P},\mathcal{P}^\prime \subseteq \operatorname{Perf}(X)$ satisfy $\overline{\langle \mathcal{P} \rangle}^{(-\infty,0]}_{\otimes} \not= \overline{\langle \mathcal{P}^\prime \rangle}^{(-\infty,0]}_{\otimes}$ and $\overline{\langle \mathbf{L} f^\ast \mathcal{P} \rangle}^{(-\infty,0]}_{\otimes} = \overline{\langle \mathbf{L} f^\ast \mathcal{P}^\prime \rangle}^{(-\infty,0]}_{\otimes}$. 
%     Then we can find an $E\in \overline{\langle \mathcal{P} \rangle}^{(-\infty,0]}_{\otimes}$ such that $E\not\in \overline{\langle \mathcal{P}^\prime \rangle}^{(-\infty,0]}_{\otimes}$. 
%     Consider the truncation triangle with respect to the aisle $\overline{\langle \mathcal{P}^\prime \rangle}^{(-\infty,0]}_{\otimes}$,
%     \begin{displaymath}
%         \tau^{\leq 0}_{\mathcal{P}^\prime} E \to E \to \tau^{\geq 1}_{\mathcal{P}^\prime}E  \to (\tau^{\leq 0}_{\mathcal{P}^\prime} E).
%     \end{displaymath}
%     By \Cref{lem:pullback_t_exactness}, the following
%     \begin{displaymath}
%         \mathbf{L} f^\ast \tau^{\leq 0}_{\mathcal{P}^\prime} E \to \mathbf{L} f^\ast E \to \mathbf{L} f^\ast \tau^{\geq 1}_{\mathcal{P}^\prime} E \to \mathbf{L} f^\ast (\tau^{\leq 0}_{\mathcal{P}^\prime} E)
%     \end{displaymath}
%     is the truncation triangle for $\mathbf{L} f^\ast E$ with respect to $\overline{\langle \mathbf{L} f^\ast \mathcal{P}^\prime \rangle}^{(-\infty,0]}_{\otimes}$. 
%     However, the assumption tells us that $\mathbf{L} f^\ast E \in \overline{\langle \mathbf{L} f^\ast \mathcal{P}^\prime \rangle}^{(-\infty,0]}_{\otimes}$, and so, $\mathbf{L} f^\ast \tau^{\geq 1}_{\mathcal{P}^\prime} E= 0$. 
%     From flatness of $f$, it follows that $\tau^{\geq 1}_{\mathcal{P}^\prime} E = 0$. 
%     This tells us that $E \in \overline{\langle \mathcal{P}^\prime \rangle}^{(-\infty,0]}_{\otimes}$, which is absurd given the assumption. 
%     Hence, we have the desired injection for \eqref{eq:classification_map1}.

%     Next we show there is an injective mapping
%     \begin{equation}
%         \label{eq:classification_map2}
%         \begin{aligned}
%             \{ \otimes\textrm{-aisle on } & D_{\operatorname{qc}}(X) \textrm{ generated by } \mathcal{P} \subseteq \operatorname{Perf}(X) \} 
%             \\&\to \{ \textrm{Thomason filtrations on } U \}
%         \end{aligned}
%     \end{equation}
%     which factors through \eqref{eq:classification_map1}. 
%     From \cite[Theorem 5.1]{Hrbek:2020}, we have an injective mapping from the image of \eqref{eq:classification_map1} into collection of Thomason filtrations on $U$. Moreover, \Cref{lem:affine_supported_by_perfects} tells us each Thomason filtration on $U$ is supported by perfects. 
%     Specifically, for each $\mathcal{P}\subseteq \operatorname{Perf}(X)$, we assign $\overline{\langle \mathbf{L} f^\ast \mathcal{P} \rangle}^{(-\infty,0]}_{\otimes}$ to the Thomason filtration on $U$ given by 
%     \begin{displaymath}
%         \phi_{\mathbf{L} f^\ast \mathcal{P}} \colon n\mapsto \bigcup_{\substack{j\geq n\\ P\in \mathcal{P}}} \operatorname{supp}(\mathcal{H}^j ( \mathbf{L}f^\ast P)).
%     \end{displaymath}
%     The presentation above follows from \Cref{lem:thomason_description}. 
%     So, tying things to together, we have the desired mapping in \eqref{eq:classification_map2}.

%     Now we claim there is a bijective mapping (i.e.\ the desired one-to-one correspondence)
%     \begin{equation}
%         \label{eq:classification_map3}
%         \begin{aligned}
%             \{ \otimes\textrm{-aisle on } & D_{\operatorname{qc}}(X) \textrm{ generated by } \mathcal{P} \subseteq \operatorname{Perf}(X) \} 
%             \\&\to \{  \textrm{Thomason filtrations on } X\}
%         \end{aligned}
%     \end{equation}
%     which factors through \eqref{eq:classification_map2}. 
%     It suffices to show there is a bijective mapping
%     \begin{equation}
%         \label{eq:classification_map4}
%         \begin{aligned}
%             \{  \textrm{Thomason filtrations on } X\} 
%             \\& \to \{ \phi_{\mathbf{L} f^\ast \mathcal{P}} \colon \mathcal{P}\subseteq \operatorname{Perf}(X) \}.
%         \end{aligned}
%     \end{equation}
%     Indeed, the image of the injective mapping in \eqref{eq:classification_map2} is precisely the target in \eqref{eq:classification_map4}. 
%     To see, by construction, the image of \eqref{eq:classification_map2} lands into $\{ \phi_{\mathbf{L} f^\ast \mathcal{P}} \colon \mathcal{P}\subseteq \operatorname{Perf}(X) \}$, so we have an injective map
%     \begin{displaymath}
%         \begin{aligned}
%             \{ \otimes\textrm{-aisle on } D_{\operatorname{qc}}(X) \textrm{ generated by } & \mathcal{P} \subseteq \operatorname{Perf}(X) \} 
%             \\&\to \{ \phi_{\mathbf{L} f^\ast \mathcal{P}} \colon \mathcal{P}\subseteq \operatorname{Perf}(X) \}.
%         \end{aligned}
%     \end{displaymath}
%     On the other hand, for each $\mathcal{P}\subseteq \operatorname{Perf}(X)$, we know $\overline{\langle \mathcal{P} \rangle}^{(-\infty,0]}_{\otimes}$ exists on $D_{\operatorname{qc}}(X)$, so its image under \eqref{eq:classification_map2} is exactly $\phi_{\mathbf{L} f^\ast \mathcal{P}}$, showing the mapping is also surjective.
    
%     Let $\psi$ be a Thomason filtration on $X$. Since it is supported by perfects, we can find $\mathcal{Q}\subseteq \operatorname{Perf}(X)$ such that 
%     \begin{displaymath}
%         \psi(n)=\bigcup_{\substack{j\geq n\\ Q\in \mathcal{Q}}} \operatorname{supp}(\mathcal{H}^j (Q)).
%     \end{displaymath}
%     Note that $g^{-1}\operatorname{supp}(R)=\operatorname{supp}(\mathbf{L}g^\ast R)$ for all morphisms $g\colon S\to T$ of quasi-compact quasi-separated algebraic spaces and $R\in \operatorname{Perf}(X)$. 
%     %%NOTE: This can be checked by taking eale presentation of target, then that of fiber product, reduce to scheme case. 
%     Using \Cref{lem:supp_perf_pullback}, it follows that 
%     \begin{displaymath}
%         \begin{aligned}
%             \bigcup_{\substack{j\geq n\\ Q\in \mathcal{Q}}} \operatorname{supp}(\mathcal{H}^j (\mathbf{L}f^\ast Q))
%             &= \bigcup_{\substack{j\geq n\\ Q\in \mathcal{Q}}} f^{-1}(\operatorname{supp}(\mathcal{H}^j (Q)))
%             \\&= f^{-1}( \bigcup_{\substack{j\geq n\\ Q\in \mathcal{Q}}} \operatorname{supp}(\mathcal{H}^j (Q))) 
%             \\&= f^{-1}(\phi(n)).
%         \end{aligned}
%     \end{displaymath}
%     It is straightforward to check that this defines a Thomason filtration on $U$. 
%     This Thomason filtration on $U$ is denoted by $f^{-1}(\phi)$. It is given by the rule $n \mapsto f^{-1}(\phi(n))$. 
%     In fact, the very construction shows $f^{-1}(\phi) = \phi_{\mathbf{L}f^\ast \mathcal{Q}}$. 
%     This gives us a well-defined mapping in \eqref{eq:classification_map4}. 
%     We need to check it is bijective. 
%     However, the construction of the mapping shows that it must be surjective, so we only need to check it is injective. 
%     Yet, this follows from the fact that $f\colon U \to X$ is surjective, completing the proof.
% \end{proof}

\bibliographystyle{alpha}
\bibliography{mainbib}

\end{document}
